% Starting a business — ICT Upper Sixth — Cours signé OnBuch+
% Official MINESEC syllabus. Compiler avec : tectonic ICT19-starting-a-business.tex
\newcommand{\DOCMATIERE}{ICT}
\newcommand{\DOCNIVEAU}{Upper Sixth}
\newcommand{\DOCMODULE}{Upper Sixth — ICT}
\newcommand{\DOCLECON}{19}
\newcommand{\DOCTITRE}{Starting a business}
% OnBuch+ — préambule « cours signés » ENRICHI (compilé avec tectonic / XeLaTeX)
% Système visuel « Pop » d'OnBuch+ : fond crème (jamais blanc), bordures encre
% épaisses, ombres « dures » décalées, titres Archivo Black en capitales.
% Version MATHÉMATIQUES : graphes (pgfplots/tikz), boîtes pédagogiques
% (définition, propriété, méthode, attention, le savais-tu, exercice résolu,
% exercice, TP, bilan), page de garde et signature OnBuch+ ancrée en bas.
%
% Macros attendues dans main.tex AVANT \input{preamble} :
%   \DOCMATIERE  \DOCNIVEAU  \DOCMODULE  \DOCLECON (numéro)  \DOCTITRE
\documentclass[11pt]{article}

\usepackage{fontspec}
\usepackage[english]{babel}
\usepackage[a4paper,margin=2cm,top=3.4cm,bottom=2.8cm,headheight=1.4cm,footskip=1.3cm]{geometry}
\usepackage{amsmath,amssymb,amsfonts}
\usepackage{mathtools} % \xmapsto, \coloneqq... souvent utilisés par les modèles
\usepackage{cancel} % simplifications barrées (bilans de forces)
% \abs{z} (module, valeur absolue) et \norm{u} (norme), utilisés par les modèles
\providecommand{\abs}{}\let\abs\relax\DeclarePairedDelimiter{\abs}{\lvert}{\rvert}
\providecommand{\norm}{}\let\norm\relax\DeclarePairedDelimiter{\norm}{\lVert}{\rVert}
\usepackage[table]{xcolor} % [table] : fournit aussi \rowcolors (alternance de lignes)
\usepackage{pagecolor}
\usepackage{tikz}
\usetikzlibrary{shapes.symbols,circuits.logic.US,circuits.logic.IEC,arrows.meta,positioning,calc,shapes.geometric,shapes.misc,shapes.arrows,decorations.pathmorphing,decorations.pathreplacing,decorations.markings,patterns,shadows,mindmap,trees,fit,backgrounds,matrix,intersections,angles,quotes}
\usepackage{pgfplots}
\usepackage[version=4]{mhchem} % équations nucléaires \ce{^{235}_{92}U}
\usepackage[european,siunitx]{circuitikz} % schémas électriques
\pgfplotsset{compat=1.18}
\usepgfplotslibrary{fillbetween,groupplots,polar,statistics,dateplot} % aires entre courbes (calcul intégral)
\usepackage{enumitem}
\usepackage{fancyhdr}
% [most] charge toutes les bibliothèques tcolorbox (listings, minted, raster,
% documentation...) et gonfle inutilement la mémoire TeX sur un document long
% (~300 boîtes) : on ne charge que ce qui est réellement utilisé.
\usepackage[skins,breakable]{tcolorbox}
\usepackage{graphicx}
\usepackage{colortbl}
\usepackage{booktabs}
\usepackage{tabularx}
\usepackage{diagbox} % case d'en-tête diagonale des tableaux à double entrée
% Colonnes extensibles centrée / gauche / droite (C, L, R), souvent utilisées par les modèles
\newcolumntype{C}{>{\centering\arraybackslash}X}
\newcolumntype{L}{>{\raggedright\arraybackslash}X}
\newcolumntype{R}{>{\raggedleft\arraybackslash}X}
\usepackage{array}
\usepackage{multicol}
\usepackage{siunitx}
% parse-numbers=false : accepte aussi \SI{2,5 \times 10^{-3}}{...}
\sisetup{locale=UK,output-decimal-marker={.},group-minimum-digits=4,parse-numbers=false}
\DeclareSIUnit\ohm{\ensuremath{\symup{\Omega}}} % U+2126 absent de la police
\DeclareSIUnit\division{div} % réglages d'oscilloscope (ms/div, V/div)
\DeclareSIUnit\year{yr}\DeclareSIUnit\annum{yr}\DeclareSIUnit\Ma{Ma}\DeclareSIUnit\tr{tr} % tours (vitesse de rotation, ex. \SI{1500}{\tr\per\minute})
\usepackage{qrcode}
% \degree : commande fréquemment utilisée par les modèles pour un angle ou une
% température en dehors de \SI{}{} (non fournie par les paquets chargés).
\providecommand{\degree}{^{\circ}}
% \qed (fin de démonstration), utilisé par les modèles sans amsthm
\providecommand{\qed}{\hfill$\square$}
% \barre : double barre des valeurs interdites dans un tableau de variations (tkz-tab)
\providecommand{\barre}{\ensuremath{\|}}
% environnement proof (amsthm non chargé) utilisé par les modèles
\ifdefined\proof\else\newenvironment{proof}[1][Proof]{\par\noindent\textit{#1.}\ }{\qed\par}\fi
\usepackage{lastpage}
\usepackage{hyperref}
\hypersetup{hidelinks}

% ============================================================
% Palette « Pop » OnBuch+ — mêmes hex que la classe P (plus_ui.dart)
% ============================================================
\definecolor{poporange}{HTML}{F59321}
\definecolor{popdark}{HTML}{E07A0C}
\definecolor{popcream}{HTML}{FFF6E8}
\definecolor{popink}{HTML}{1C1714}
\definecolor{poppurple}{HTML}{7A5AE0}
\definecolor{popgreen}{HTML}{2E9E6A}
\definecolor{poppink}{HTML}{E0457B}
\definecolor{popblue}{HTML}{2F7DE1}
\definecolor{popgold}{HTML}{C98A12}
\definecolor{popmuted}{HTML}{8A7F76}
\definecolor{popline}{HTML}{EADFD2}
\definecolor{popgray}{HTML}{8A7F76}
\colorlet{popgrey}{popgray}
% Alias tolérants (le modèle nomme parfois les couleurs différemment)
\colorlet{popwhite}{white}
\colorlet{popblack}{popink}
\colorlet{popred}{poppink}
\colorlet{popyellow}{popgold}
% Teintes claires pour fonds de tableaux / zones
\colorlet{popblueL}{popblue!10!white}
\colorlet{poporangeL}{poporange!14!white}
\colorlet{popgreenL}{popgreen!12!white}
\colorlet{poppurpleL}{poppurple!12!white}
\colorlet{poppinkL}{poppink!12!white}
\colorlet{popgoldL}{popgold!14!white}

\pagecolor{popcream}

% ============================================================
% Typographie
% ============================================================
\defaultfontfeatures{Ligatures=TeX}
\setmainfont{PlusJakartaSans-Medium.ttf}[Path=../../../fonts/,BoldFont=PlusJakartaSans-Bold.ttf]
% Maths en sans-serif (Fira Math) avec chiffres et lettres droites de Plus
% Jakarta, pour que les formules se fondent dans le texte.
\usepackage[math-style=ISO,bold-style=ISO]{unicode-math}
\setmathfont{FiraMath-Regular.otf}[Path=../../../fonts/]
\setmathfont{PlusJakartaSans-Medium.ttf}[Path=../../../fonts/,range={up/num,up/{Latin,latin},\mathup}]
% (icomma retiré : décimales avec point en anglais)
% Caractères mathématiques Unicode absents des polices de texte (Plus Jakarta,
% Archivo Black) : rendus par leur équivalent mathématique, en texte comme en
% formule — sinon ils s'affichent en carré vide (ex. « Arithmétique dans ℤ »).
\usepackage{newunicodechar}
\newunicodechar{ℝ}{\ensuremath{\mathbb{R}}}
\newunicodechar{ℤ}{\ensuremath{\mathbb{Z}}}
\newunicodechar{ℂ}{\ensuremath{\mathbb{C}}}
\newunicodechar{ℕ}{\ensuremath{\mathbb{N}}}
\newunicodechar{ℚ}{\ensuremath{\mathbb{Q}}}
\newunicodechar{𝔻}{\ensuremath{\mathbb{D}}}
\newunicodechar{α}{\ensuremath{\alpha}}
\newunicodechar{β}{\ensuremath{\beta}}
\newunicodechar{γ}{\ensuremath{\gamma}}
\newunicodechar{δ}{\ensuremath{\delta}}
\newunicodechar{ε}{\ensuremath{\varepsilon}}
\newunicodechar{θ}{\ensuremath{\theta}}
\newunicodechar{λ}{\ensuremath{\lambda}}
\newunicodechar{μ}{\ensuremath{\mu}}
\newunicodechar{σ}{\ensuremath{\sigma}}
\newunicodechar{τ}{\ensuremath{\tau}}
\newunicodechar{φ}{\ensuremath{\varphi}}
\newunicodechar{ψ}{\ensuremath{\psi}}
\newunicodechar{ω}{\ensuremath{\omega}}
\newunicodechar{Σ}{\ensuremath{\Sigma}}
% majuscules produites par \MakeUppercase dans les titres (α-aminés -> Α)
\newunicodechar{Α}{\ensuremath{\alpha}}
\newunicodechar{Β}{\ensuremath{\beta}}
\newunicodechar{Γ}{\ensuremath{\Gamma}}
\newunicodechar{Δ}{\ensuremath{\Delta}}
\newunicodechar{Ε}{\ensuremath{\varepsilon}}
\newunicodechar{Θ}{\ensuremath{\Theta}}
\newunicodechar{Λ}{\ensuremath{\Lambda}}
\newunicodechar{Μ}{\ensuremath{\mu}}
\newunicodechar{Τ}{\ensuremath{\tau}}
\newunicodechar{Φ}{\ensuremath{\Phi}}
\newunicodechar{Ψ}{\ensuremath{\Psi}}
\newunicodechar{Ω}{\ensuremath{\Omega}}
\newunicodechar{↦}{\ensuremath{\mapsto}}
\newunicodechar{∈}{\ensuremath{\in}}
\newunicodechar{∉}{\ensuremath{\notin}}
\newunicodechar{∧}{\ensuremath{\wedge}}
\newunicodechar{⇔}{\ensuremath{\Leftrightarrow}}
\newunicodechar{⟺}{\ensuremath{\Longleftrightarrow}}
\newunicodechar{⇒}{\ensuremath{\Rightarrow}}
\newunicodechar{⟹}{\ensuremath{\Longrightarrow}}
\newunicodechar{∘}{\ensuremath{\circ}}
\newunicodechar{∪}{\ensuremath{\cup}}
\newunicodechar{∩}{\ensuremath{\cap}}
\newunicodechar{≡}{\ensuremath{\equiv}}
\newunicodechar{⊂}{\ensuremath{\subset}}
\newunicodechar{⊥}{\ensuremath{\perp}}
\newunicodechar{∃}{\ensuremath{\exists}}
\newunicodechar{∀}{\ensuremath{\forall}}
\newunicodechar{∛}{\ensuremath{\sqrt[3]{\,}}}
\newunicodechar{∜}{\ensuremath{\sqrt[4]{\,}}}
\newunicodechar{⃗}{\ensuremath{{}^{\to}}}
\newunicodechar{ⁿ}{\ifmmode^{n}\else\textsuperscript{\ensuremath{n}}\fi}
\newunicodechar{ˣ}{\ifmmode^{x}\else\textsuperscript{\ensuremath{x}}\fi}
\newunicodechar{ᵇ}{\ifmmode^{b}\else\textsuperscript{\ensuremath{b}}\fi}
\newunicodechar{ʳ}{\ifmmode^{r}\else\textsuperscript{\ensuremath{r}}\fi}
\newunicodechar{ᵘ}{\ifmmode^{u}\else\textsuperscript{\ensuremath{u}}\fi}
\newunicodechar{ʸ}{\ifmmode^{y}\else\textsuperscript{\ensuremath{y}}\fi}
\newunicodechar{ᵃ}{\ifmmode^{a}\else\textsuperscript{\ensuremath{a}}\fi}
\newunicodechar{ᵉ}{\ifmmode^{e}\else\textsuperscript{\ensuremath{e}}\fi}
\newunicodechar{ᵏ}{\ifmmode^{k}\else\textsuperscript{\ensuremath{k}}\fi}
\newunicodechar{ᵖ}{\ifmmode^{p}\else\textsuperscript{\ensuremath{p}}\fi}
\newunicodechar{ᴺ}{\ifmmode^{N}\else\textsuperscript{\ensuremath{N}}\fi}
\newunicodechar{ᶜ}{\ifmmode^{c}\else\textsuperscript{\ensuremath{c}}\fi}
\newunicodechar{ʰ}{\ifmmode^{h}\else\textsuperscript{\ensuremath{h}}\fi}
\newunicodechar{ᵠ}{\ifmmode^{q}\else\textsuperscript{\ensuremath{q}}\fi}
\newunicodechar{ₙ}{\ifmmode_{n}\else\textsubscript{\ensuremath{n}}\fi}
\newunicodechar{ₐ}{\ifmmode_{a}\else\textsubscript{\ensuremath{a}}\fi}
\newunicodechar{ᵢ}{\ifmmode_{i}\else\textsubscript{\ensuremath{i}}\fi}
\newunicodechar{ᵣ}{\ifmmode_{r}\else\textsubscript{\ensuremath{r}}\fi}
\newunicodechar{ₖ}{\ifmmode_{k}\else\textsubscript{\ensuremath{k}}\fi}
\newunicodechar{ᵦ}{\ifmmode_{\beta}\else\textsubscript{\ensuremath{\beta}}\fi}
\newunicodechar{ₓ}{\ifmmode_{x}\else\textsubscript{\ensuremath{x}}\fi}
\newfontfamily\popdisplay{ArchivoBlack-Regular.ttf}[Path=../../../fonts/]
\newfontfamily\poplogo{SpaceGrotesk-Bold.ttf}[Path=../../../fonts/]
\newfontfamily\popsemi{PlusJakartaSans-SemiBold.ttf}[Path=../../../fonts/]
\newfontfamily\popxbold{PlusJakartaSans-ExtraBold.ttf}[Path=../../../fonts/]
\newfontfamily\popxboldls{PlusJakartaSans-ExtraBold.ttf}[Path=../../../fonts/,LetterSpace=3.0]

\setlist[enumerate]{leftmargin=1.7em,itemsep=5pt,topsep=4pt}
\setlist[itemize]{leftmargin=1.7em,itemsep=4pt,topsep=4pt,label={\color{poporange}\textbullet}}
\setlength{\parindent}{0pt}
\setlength{\parskip}{5pt}
\renewcommand{\arraystretch}{1.25}

% pgfplots : style Pop par défaut
\pgfplotsset{
  every axis/.append style={axis line style={popink,line width=1pt},tick style={popink},
    grid style={popline,line width=0.5pt},label style={font=\small},tick label style={font=\footnotesize}},
  popcurve/.style={poporange,line width=1.8pt,no markers,smooth},
}
% Même style disponible dans un \draw TikZ simple (hors pgfplots)
\tikzset{popcurve/.style={poporange,line width=1.8pt}}
\tikzset{popgrid/.style={popline,very thin}}
\colorlet{popcurve}{poporange} % le nom du style est parfois utilisé comme couleur

% ============================================================
% Pill / squiggle
% ============================================================
\newcommand{\poppill}[3]{%
  \tikz[baseline=-0.55ex]\node[draw=popink,line width=1.2pt,rounded corners=2.6pt,%
    fill=#2,inner xsep=6.5pt,inner ysep=3pt]{%
    \popxboldls\fontsize{7}{7}\selectfont\color{#3}\MakeUppercase{#1}};%
}
\newcommand{\popsquiggle}[1]{%
  \tikz[baseline=-0.4ex]{
    \draw[#1,line width=2.6pt,line cap=round]
      plot [smooth] coordinates {(0,0.09) (0.34,0.01) (0.68,0.21) (1.02,0.01) (1.36,0.10)};
  }%
}
% Mot-clé surligné (marqueur orange sous le texte)
\newcommand{\cle}[1]{{\popxbold\color{popdark}#1}}

% ============================================================
% En-tête / pied
% ============================================================
\pagestyle{fancy}
\fancyhf{}
\renewcommand{\headrulewidth}{0pt}
\renewcommand{\footrulewidth}{0pt}
\lhead{\raisebox{-0.15ex}{\poplogo\fontsize{13}{13}\selectfont\color{popink}OnBuch\color{poporange}+}}
\rhead{\poppill{\DOCMATIERE\ \textbullet\ \DOCNIVEAU\ \textbullet\ Chapter \DOCLECON}{white}{popink}}
\lfoot{\popxbold\fontsize{7.6}{7.6}\selectfont\color{popmuted}\MakeUppercase{Signed course OnBuch+}\ \textbullet\ {\popsemi\DOCTITRE}}
\rfoot{\tikz[baseline=-0.6ex]\node[rounded corners=8.5pt,fill=popink,minimum height=17pt,minimum width=17pt,inner xsep=5pt,inner sep=0]{\popxbold\fontsize{8}{8}\selectfont\color{popcream}\thepage\,/\,\pageref*{LastPage}};}

% ============================================================
% Titres
% ============================================================
\newcommand{\chaptitre}[2]{%
  \begin{tcolorbox}[enhanced,colback=poporange,colframe=popink,boxrule=2.6pt,arc=11pt,
      drop shadow=popink,left=15pt,right=15pt,top=12pt,bottom=11pt,before skip=3pt,after skip=18pt]
    \poppill{#2}{popink}{popcream}\par\vspace{9pt}
    {\raggedright\hyphenpenalty=10000\exhyphenpenalty=10000\popdisplay\fontsize{22}{24}\selectfont\color{popcream}\MakeUppercase{#1}\par}
  \end{tcolorbox}
}
% Section numérotée : pastille orange avec le numéro + titre Archivo
\newcounter{popsec}
\newcommand{\coursec}[1]{%
  \stepcounter{popsec}\setcounter{popsub}{0}%
  \par\vspace{16pt}\noindent
  \begin{minipage}{\linewidth}
  \tikz[baseline=-0.7ex]\node[draw=popink,line width=1.6pt,fill=poporange,rounded corners=4pt,minimum size=22pt,inner sep=0]{\popdisplay\fontsize{12}{12}\selectfont\color{popcream}\thepopsec};%
  \hspace{8pt}{\popdisplay\fontsize{15}{17}\selectfont\color{popink}\MakeUppercase{#1}}\par
  \vspace{4pt}\noindent{\color{poporange}\rule{36pt}{3.6pt}}
  \end{minipage}\par\nopagebreak\vspace{8pt}
}
\newcounter{popsub}
\newcommand{\courssub}[1]{%
  \stepcounter{popsub}%
  \par\vspace{9pt}\noindent{\popxbold\fontsize{11.5}{13}\selectfont\color{popink}{\color{poporange}\thepopsec.\thepopsub}\hspace{6pt}#1}\par\nopagebreak\vspace{4pt}
}

% ============================================================
% Boîtes pédagogiques — même recette : fond blanc, bordure épaisse colorée,
% ombre dure encre, badge plein. Titre optionnel [#1].
% ============================================================
\tcbset{popbase/.style={breakable,enhanced,colback=white,boxrule=2.3pt,arc=9pt,
  shadow={3pt}{-3pt}{0pt}{popink},left=14pt,right=14pt,top=11pt,bottom=11pt,before skip=11pt,after skip=15pt,
  fontupper=\popsemi\fontsize{10.6}{14.5}\selectfont\color{popink},
  fontlower=\popsemi\fontsize{10.6}{14.5}\selectfont\color{popink}}}
% Coche dessinée (le glyphe ✓ n'existe pas dans les polices du document)
\newcommand{\popcheck}[1]{\tikz[baseline=-0.5ex]\draw[#1,line width=1.6pt,line cap=round,line join=round](-0.13,0.0)--(-0.04,-0.09)--(0.13,0.1);}
\providecommand{\coche}{\popcheck{popgreen}}
\providecommand{\resultat}[1]{\par\smallskip\noindent\fcolorbox{popgreen}{popgreenL}{\parbox{\dimexpr\linewidth-2\fboxsep-2\fboxrule\relax}{\textbf{Result.} #1}}\par\smallskip}
\newcommand{\popboxhead}[3]{\poppill{#1}{#2}{white}\ifx\relax#3\relax\else\hspace{6pt}{\popxbold\fontsize{10.6}{12}\selectfont\color{popink}#3}\fi\par\vspace{7pt}}

\newtcolorbox{aretenir}[1][]{popbase,colframe=popgreen,before upper={\popboxhead{Key points}{popgreen}{#1}}}
\newtcolorbox{exemplebox}[1][]{popbase,colframe=poporange,before upper={\popboxhead{Exemple}{poporange}{#1}}}
\newtcolorbox{definition}[1][]{popbase,colframe=popblue,before upper={\popboxhead{Definition}{popblue}{#1}}}
\newtcolorbox{propriete}[1][]{popbase,colframe=poppurple,before upper={\popboxhead{Property}{poppurple}{#1}}}
\newtcolorbox{methode}[1][]{popbase,colframe=popgold,before upper={\popboxhead{Method}{popgold}{#1}}}
\newtcolorbox{attention}[1][]{popbase,colframe=poppink,before upper={\popboxhead{Warning}{poppink}{#1}}}
\newtcolorbox{experience}[1][]{popbase,colframe=popblue,colback=popblueL,before upper={\popboxhead{Experiment}{popblue}{#1}}}
% « Le savais-tu ? » : carte violette pleine (contexte, culture, Cameroun)
\newtcolorbox{savaistu}[1][]{popbase,colback=poppurple,colframe=popink,
  fontupper=\popsemi\fontsize{10.4}{14.2}\selectfont\color{white},
  before upper={\poppill{Did you know?}{white}{poppurple}\ifx\relax#1\relax\else\hspace{6pt}{\popxbold\color{white}#1}\fi\par\vspace{7pt}}}
% Exercice résolu : énoncé en haut, \tcblower puis la solution sur fond crème
\newcounter{popexo}\newcounter{popexr}
\newtcolorbox{exoresolu}[1][]{popbase,colframe=poporange,colbacklower=popcream,
  segmentation style={popink,line width=1.2pt,dashed},
  before upper={\stepcounter{popexr}\popboxhead{Worked example \thepopexr}{poporange}{#1}},
  before lower={\poppill{Solution}{popink}{popcream}\par\vspace{6pt}}}
% Exercice (non résolu en place) — la correction est dans la partie Corrigés
\newtcolorbox{exercice}[1][]{popbase,colframe=popink,
  before upper={\stepcounter{popexo}\popboxhead{Exercise \thepopexo}{popink}{#1}}}
% Corrigé
\newtcolorbox{corrige}[1][]{popbase,colframe=popgreen,colback=popgreenL,
  before upper={\popboxhead{Solution}{popgreen}{#1}}}
% Objectifs / prérequis / bilan
\newtcolorbox{objectifs}{popbase,colframe=popink,colback=poporangeL,
  before upper={\popboxhead{Chapter objectives}{popink}{}}}
\newtcolorbox{prerequis}{popbase,colframe=popink,colback=popblueL,
  before upper={\popboxhead{Prerequisites}{popblue}{}}}
\newtcolorbox{bilan}{popbase,colframe=popink,colback=white,boxrule=2.8pt,
  before upper={\popboxhead{Fiche bilan}{poporange}{L'essentiel à connaître pour l'examen}}}
% Figure encadrée avec légende
\newtcolorbox{popfigure}[1][]{enhanced,breakable=false,colback=white,colframe=popline,boxrule=1.4pt,arc=8pt,
  left=10pt,right=10pt,top=10pt,bottom=8pt,before skip=10pt,after skip=12pt,halign=center,
  lower separated=false,halign lower=center,
  fontlower=\popsemi\fontsize{9}{11.5}\selectfont\color{popmuted},
  #1}
\newcommand{\legende}[1]{\par\vspace{6pt}{\popsemi\fontsize{9}{11.5}\selectfont\color{popmuted}#1}}

% ============================================================
% Signature OnBuch+
% ============================================================
\newcommand{\poplogoinline}[1]{{\poplogo\fontsize{#1}{#1}\selectfont\color{popink}OnBuch\color{poporange}+}}
\newcommand{\signatureonbuch}{%
  \begin{tcolorbox}[enhanced,colback=popink,colframe=popink,boxrule=2.2pt,arc=10pt,
      drop shadow=poporange,left=14pt,right=14pt,top=10pt,bottom=10pt,before skip=0pt,after skip=0pt]
    \tikz[baseline=-0.7ex]\node[circle,fill=popgreen,draw=popcream,line width=1.3pt,minimum size=22pt,inner sep=0]{\popcheck{white}};%
    \hspace{9pt}\begin{minipage}[c]{0.62\linewidth}
      {\popxbold\fontsize{12}{13}\selectfont\color{popcream}Signed\ \textbullet\ {\poplogo OnBuch\color{poporange}+}}\par\vspace{2pt}
      {\raggedright\popsemi\fontsize{8}{10.5}\selectfont\color{popline}\MakeUppercase{OnBuch+ teaching team}\\\MakeUppercase{Follows the official MINESEC syllabus}\par}
    \end{minipage}\hfill
    {\poplogo\fontsize{22}{22}\selectfont\color{popcream}OnBuch\color{poporange}+}
  \end{tcolorbox}%
}

% ============================================================
% Page de garde
% #1 titre · #2 matière · #3 niveau · #4 module · #5 n° leçon · #6 durée ·
% #7 description / objectifs (texte ou itemize)
% ============================================================
\newcommand{\pagegarde}[7]{%
  \thispagestyle{empty}
  \newgeometry{a4paper,margin=1.7cm,top=1.6cm,bottom=1.4cm}%
  \begin{tikzpicture}[remember picture,overlay]
    \fill[poporange!18] ([xshift=-3.2cm,yshift=-3.6cm]current page.north east) circle (4.6cm);
    \fill[poppurple!14] ([xshift=2.4cm,yshift=7.6cm]current page.south west) circle (3.8cm);
  \end{tikzpicture}%
  {\poplogo\fontsize{30}{30}\selectfont\color{popink}OnBuch\color{poporange}+}\hfill
  \raisebox{0.6ex}{\poppill{Signed course \textbullet\ Premium}{popink}{popcream}}\par
  \vspace{1pt}\popsquiggle{poppurple}\par
  \vspace{8pt}
  \begin{tcolorbox}[enhanced,colback=poporange,colframe=popink,boxrule=2.8pt,arc=13pt,
      drop shadow=popink,left=18pt,right=18pt,top=12pt,bottom=13pt,after skip=10pt]
    \poppill{#2 \textbullet\ #3}{popink}{popcream}\hspace{5pt}\poppill{#4}{white}{popink}\par\vspace{8pt}
    {\popdisplay\fontsize{14}{14}\selectfont\color{popink}CHAPTER #5}\par\vspace{4pt}
    {\raggedright\hyphenpenalty=10000\exhyphenpenalty=10000\popdisplay\fontsize{25}{27}\selectfont\color{popcream}\MakeUppercase{#1}\par}
    \vspace{8pt}
    {\popxbold\fontsize{9.5}{9.5}\selectfont\color{popink}\MakeUppercase{Suggested time: #6}}
  \end{tcolorbox}
  \begin{tcolorbox}[enhanced,colback=white,colframe=popink,boxrule=2.2pt,arc=10pt,
      drop shadow=popink,left=16pt,right=16pt,top=10pt,bottom=10pt,after skip=8pt]
    \poppill{In this chapter}{poppurple}{white}\par\vspace{5pt}
    \popsemi\fontsize{9.3}{12.6}\selectfont\color{popink}#7
  \end{tcolorbox}
  \noindent\begin{minipage}[t]{0.487\linewidth}
    \begin{tcolorbox}[enhanced,colback=popgreen,colframe=popink,boxrule=2.2pt,arc=10pt,
        drop shadow=popink,left=11pt,right=11pt,top=10pt,bottom=10pt,height=3.1cm,valign=center]
      \tcbox[colback=white,colframe=popink,boxrule=1.3pt,arc=5pt,boxsep=0pt,nobeforeafter,
        left=3pt,right=3pt,top=3pt,bottom=3pt,valign=center]{\qrcode[height=1.9cm]{https://play.google.com/store/apps/details?id=com.luvvix.onbuch&hl=en}}%
      \hspace{8pt}\begin{minipage}[c]{0.5\linewidth}
        {\popdisplay\fontsize{11}{12.5}\selectfont\color{white}\MakeUppercase{Download OnBuch}}\par\vspace{4pt}
        {\raggedright\popsemi\fontsize{8.4}{11}\selectfont\color{white}Free \textbullet\ Google Play\\AI tutor, past papers, results\par}
      \end{minipage}
    \end{tcolorbox}
  \end{minipage}\hfill
  \begin{minipage}[t]{0.487\linewidth}
    \begin{tcolorbox}[enhanced,colback=poppurple,colframe=popink,boxrule=2.2pt,arc=10pt,
        drop shadow=popink,left=11pt,right=11pt,top=10pt,bottom=10pt,height=3.1cm,valign=center]
      \tikz[baseline=-0.7ex]\node[circle,fill=white,draw=popink,line width=1.3pt,minimum size=1.6cm,inner sep=0]{\popdisplay\fontsize{24}{24}\selectfont\color{poppurple}+};%
      \hspace{8pt}\begin{minipage}[c]{0.6\linewidth}
        {\popdisplay\fontsize{11}{12.5}\selectfont\color{white}\MakeUppercase{Go OnBuch+}}\par\vspace{4pt}
        {\raggedright\popsemi\fontsize{8.4}{11}\selectfont\color{white}All signed courses, unlimited AI tutor, PDF sheets — OnBuch+ tab in the app\par}
      \end{minipage}
    \end{tcolorbox}
  \end{minipage}
  \vspace{8pt}
  \popcontenu
  \vfill
  \signatureonbuch
  \restoregeometry
  \newpage
}

% Rangée « Ce cours contient » (page de garde)
\newcommand{\poptile}[3]{% #1 couleur #2 gros texte #3 légende
  \begin{tcolorbox}[enhanced,colback=#1,colframe=popink,boxrule=2pt,arc=9pt,shadow={2.5pt}{-2.5pt}{0pt}{popink},
      left=6pt,right=6pt,top=9pt,bottom=9pt,halign=center,valign=center,height=1.75cm,width=\linewidth,nobeforeafter]
    {\popdisplay\fontsize{12.5}{13}\selectfont\color{white}\MakeUppercase{#2}}\par\vspace{3pt}
    {\centering\popsemi\fontsize{7.8}{9.5}\selectfont\color{white}#3\par}
  \end{tcolorbox}}
\newcommand{\popcontenu}{%
  \par\noindent{\popxboldls\fontsize{7.5}{7.5}\selectfont\color{popmuted}THIS SIGNED COURSE CONTAINS}\par\vspace{7pt}
  \noindent\begin{minipage}[t]{0.235\linewidth}\poptile{popblue}{Course}{complete, illustrated, step by step}\end{minipage}\hfill
  \begin{minipage}[t]{0.235\linewidth}\poptile{poporange}{Methods}{and worked examples}\end{minipage}\hfill
  \begin{minipage}[t]{0.235\linewidth}\poptile{poppink}{Exercises}{exam style + solutions}\end{minipage}\hfill
  \begin{minipage}[t]{0.235\linewidth}\poptile{popgreen}{Summary sheet}{the essentials to remember}\end{minipage}\par}

% Sommaire
\newtcolorbox{sommaire}{enhanced,colback=white,colframe=popink,boxrule=2.2pt,arc=10pt,
  drop shadow=popink,left=18pt,right=18pt,top=13pt,bottom=13pt,before skip=6pt,after skip=20pt,
  before upper={\poppill{Contents}{poppurple}{white}\par\vspace{9pt}\popsemi\fontsize{10.6}{14.5}\selectfont\color{popink}}}

% Signature de fin de cours — poussée tout en bas de la dernière page
\newcommand{\findecours}{%
  \par\vfill\nopagebreak
  \begin{center}\popsquiggle{poporange}\end{center}\vspace{4pt}
  \signatureonbuch
}
\AtBeginDocument{\def\llbracket{\mathopen{[\mkern-3.5mu[}}\def\rrbracket{\mathclose{]\mkern-3.5mu]}}} % intervalles d'entiers (glyphe ⟦ absent de la police)
% Glyphes absents de Fira Math : redéfinis à partir de glyphes disponibles
\AtBeginDocument{%
  \def\setminus{\mathbin{\backslash}}\let\smallsetminus\setminus
  \def\vdots{\mathord{\vbox{\baselineskip3.2pt\lineskiplimit0pt\kern4pt\hbox{.}\hbox{.}\hbox{.}}}}%
  \def\ddots{\mathinner{\mkern1mu\raise6.5pt\hbox{.}\mkern2mu\raise3.6pt\hbox{.}\mkern2mu\raise0.7pt\hbox{.}\mkern1mu}}%
  \def\triangle{\mathord{\scalebox{0.9}{$\symup{\Delta}$}}}%
  \def\nabla{\mathord{\raisebox{\depth}{\scalebox{1}[-1]{$\symup{\Delta}$}}}}%
  \def\checkmark{\popcheck{popdark}}%
  \def\star{\mathbin{\ast}}\def\bigstar{\mathord{\ast}}%
  \def\diamond{\mathbin{\scalebox{0.6}{$\Diamond$}}}%
  \def\bigcup{\mathop{\mathchoice{\raisebox{-0.3em}{\scalebox{1.7}{$\cup$}}}{\scalebox{1.25}{$\cup$}}{\cup}{\cup}}\displaylimits}%
  \def\bigcap{\mathop{\mathchoice{\raisebox{-0.3em}{\scalebox{1.7}{$\cap$}}}{\scalebox{1.25}{$\cap$}}{\cap}{\cap}}\displaylimits}%
  \def\lesseqqgtr{\mathrel{\lessgtr}}\def\bigoplus{\mathop{\oplus}\displaylimits}%
}
\providecommand{\ssi}{if and only if} % abréviation fréquente
\AtBeginDocument{\providecommand{\wideparen}{\overparen}} % arc de cercle

% Fonctions usuelles que les modèles utilisent sans que amsmath les fournisse
\providecommand{\arsinh}{\operatorname{arsinh}}\providecommand{\arcosh}{\operatorname{arcosh}}
\providecommand{\artanh}{\operatorname{artanh}}\providecommand{\arsech}{\operatorname{arsech}}
\providecommand{\arcsch}{\operatorname{arcsch}}\providecommand{\arcoth}{\operatorname{arcoth}}
\providecommand{\arcsinh}{\operatorname{arsinh}}\providecommand{\arccosh}{\operatorname{arcosh}}
\providecommand{\arctanh}{\operatorname{artanh}}\providecommand{\sech}{\operatorname{sech}}
\providecommand{\csch}{\operatorname{csch}}\providecommand{\arcsec}{\operatorname{arcsec}}
\providecommand{\arccsc}{\operatorname{arccsc}}\providecommand{\arccot}{\operatorname{arccot}}
\providecommand{\sgn}{\operatorname{sgn}}\providecommand{\lcm}{\operatorname{lcm}}
\providecommand{\Var}{\operatorname{Var}}\providecommand{\Cov}{\operatorname{Cov}}
\providecommand{\permil}{\ensuremath{{}^{0}\!/_{00}}}\providecommand{\textperthousand}{\ensuremath{{}^{0}\!/_{00}}}

%% FIGFIX-BEGIN v4 — correctifs globaux des figures (docs/onbuch-generation/tools/figfix)
\usepackage{adjustbox}\usepackage{etoolbox}
% toute figure TikZ/pgfplots plus large que la ligne est réduite à la largeur disponible
\BeforeBeginEnvironment{tikzpicture}{\begin{adjustbox}{max width=\linewidth}}
\AfterEndEnvironment{tikzpicture}{\end{adjustbox}}
% légendes pgfplots lisibles par-dessus les courbes
\pgfplotsset{every axis/.append style={legend style={fill=white,fill opacity=0.92,text opacity=1,draw=black!15,inner sep=3pt}}}
% étiquettes d'arêtes (syntaxe edge["..."]) avec fond blanc
\tikzset{every edge quotes/.append style={fill=white,inner sep=1.2pt}}
%% FIGFIX-END
\begin{document}

\pagegarde{Starting a business}{ICT}{Upper Sixth}{Upper Sixth — ICT}{19}{1 periods}{This chapter introduces students to the first step of IT-based entrepreneurship: observing a business environment, identifying real needs and problems, and determining which of them can be solved with information technology. Through Cameroonian case studies, learners develop the analytical skills required to propose relevant, feasible IT solutions — a key competence for the GCE Advanced Level and for the world of work.
\vspace{4pt}
\begin{multicols}{2}\small
\begin{itemize}
\item By the end of this chapter I can define the concepts of business, business need and business opportunity.
\item By the end of this chapter I can distinguish between needs that are solvable with IT and those that are not.
\item By the end of this chapter I can observe a real business environment and collect information about its problems using surveys, interviews and observation.
\item By the end of this chapter I can classify identified business needs by domain (communication, record keeping, marketing, payments, security, etc.).
\item By the end of this chapter I can match common business problems to appropriate IT tools and services.
\item By the end of this chapter I can evaluate the feasibility of an IT solution in terms of cost, skills, infrastructure and connectivity in the Cameroonian context.
\end{itemize}\end{multicols}}

\begin{sommaire}
\begin{multicols}{2}
\begin{enumerate}
\item Understanding Businesses and Their Needs
\item Techniques for Identifying Business Needs
\item Matching Business Needs to IT Solutions
\item Evaluating Feasibility and Spotting IT Business Opportunities
\item Integration activity
\item Methods and worked examples
\item Exercises
\item Solutions to the exercises
\item Summary sheet
\end{enumerate}
\end{multicols}
\end{sommaire}

\begin{objectifs}
\begin{itemize}
\item By the end of this chapter I can define the concepts of business, business need and business opportunity.
\item By the end of this chapter I can distinguish between needs that are solvable with IT and those that are not.
\item By the end of this chapter I can observe a real business environment and collect information about its problems using surveys, interviews and observation.
\item By the end of this chapter I can classify identified business needs by domain (communication, record keeping, marketing, payments, security, etc.).
\item By the end of this chapter I can match common business problems to appropriate IT tools and services.
\item By the end of this chapter I can evaluate the feasibility of an IT solution in terms of cost, skills, infrastructure and connectivity in the Cameroonian context.
\item By the end of this chapter I can write a simple problem statement and a needs analysis report for a small business.
\item By the end of this chapter I can identify business opportunities created by IT in Cameroon (e-commerce, mobile money services, digital agriculture, online education).
\end{itemize}
\end{objectifs}

\begin{prerequis}
\begin{itemize}
\item Basic computer hardware and software concepts (Lower Sixth)
\item Common application software: word processors, spreadsheets, databases
\item Internet services: email, websites, social networks, search engines
\item Mobile telephony and Mobile Money usage in everyday Cameroonian life
\item Notion of information systems and data processing
\end{itemize}
\end{prerequis}

\newpage

\coursec{Understanding Businesses and Their Needs}

Every Saturday morning, Mokolo Market in Yaoundé wakes up before dawn. Among the hundreds of stalls, Mama Ngo Bell sells tomatoes, onions and peppers. She is hardworking, her prices are fair, and her tomatoes are fresh. Yet every Saturday, by 11 a.m., her stall is empty of tomatoes while customers are still queuing. She has lost count of the number of buyers who walked away to her competitors. Her problem is simple: she cannot track her stock. She buys ``by feeling'', records nothing, and only discovers the shortage when it is too late.

Next door, a young trader named Sandrine sells almost the same products. But Sandrine uses a simple smartphone application to record every sale, checks which days sell the most, predicts her demand for the coming week, and even sends Mobile Money payment links to customers who order in advance. Her stall is always stocked, her customers rarely leave disappointed, and her profits grow month after month.

Same market. Same products. Same customers. The only difference is the \cle{use of information technology} to answer a business need. What if the difference between a struggling business and a thriving one was simply the ability to \emph{spot a problem that IT can solve}? This chapter trains you to see business needs the way an IT entrepreneur does. Before we can identify such needs (Lesson 51 of the official syllabus), we must first understand what a business is, what a business need is, and why information is the raw material of good management.

\courssub{What Is a Business?}

\textbf{Intuition first.} Look around any Cameroonian town: the call box operator at the junction, the bakery that supplies your quarter with bread, the provision store near the school, the transport agency on the Douala--Yaoundé highway, the cybercafé, the mobile money kiosk. All these activities look different, yet they share the same deep structure: each one \emph{provides something} (a good or a service) \emph{to someone} (a customer) \emph{in exchange for money}, and each one is run with a \emph{goal} — usually profit, sometimes a social mission.

\begin{definition}[Business]
A \cle{business} is an organisation or activity that provides \textbf{goods} (physical products such as bread, tomatoes or phone credit) or \textbf{services} (work done for others, such as transport, hair dressing or internet access) to \textbf{customers} in exchange for money, with the aim of making a \textbf{profit} or achieving a \textbf{social goal}.
\end{definition}

Three elements of this definition deserve attention. First, a business always involves an \emph{exchange}: value flows to the customer, money flows back. Second, a business has \emph{customers} — people or organisations whose needs it satisfies. Third, a business has a \emph{purpose}: most private businesses seek profit (the difference between revenue and costs), while cooperatives and social enterprises may primarily pursue a community goal, such as helping farmers sell their cocoa at a fair price.

\textbf{Types of businesses in Cameroon.} The Cameroonian economic landscape contains several legal and practical forms of business, which you must be able to recognise:

\begin{center}
\begin{tabularx}{\linewidth}{|l|X|X|}
\hline
\rowcolor{popblueL} \textbf{Type} & \textbf{Description} & \textbf{Cameroonian examples} \\
\hline
Sole proprietorship & Owned and run by one person; simple to create; the owner keeps all profits but bears all risks. & Call boxes, provision stores, tailoring workshops, ``bend-skin'' repair points. \\
\hline
Partnership & Two or more persons pool money and skills and share profits and responsibilities. & Two nurses opening a private clinic; friends running a printing press. \\
\hline
Company (SARL, SA) & A legal entity separate from its owners. A \textbf{SARL} (limited liability company) suits small/medium firms; an \textbf{SA} (public limited company) suits large ones that can raise capital from many shareholders. & A SARL distributing beverages in Bamenda; large SAs in banking, insurance or telecommunications. \\
\hline
Cooperative / CIG & Members unite voluntarily to meet common economic needs; decisions are democratic. & Common Initiative Groups (CIGs) of farmers in the North West; cocoa or coffee cooperatives; women's savings groups (``njangi'' formalised). \\
\hline
Informal sector & Activities not formally registered, often small-scale and family-based, but economically vital. & Hawkers, roadside food sellers (``tournedos''), motorcycle transporters before formal registration. \\
\hline
\end{tabularx}
\end{center}

\begin{attention}[Size and formality do not change the definition]
A one-woman tomato stall and a national telecommunications company are \emph{both} businesses. Both provide value to customers, both receive money, both must manage stock, money, records and customers. Consequently, \textbf{both have business needs that IT can address} — only the scale of the solution differs.
\end{attention}

\courssub{What Is a Business Need?}

\textbf{Intuition first.} A human being needs air, water, food and shelter to survive, and needs education and work to grow. A business is similar: it needs certain things to \emph{operate} day to day, to \emph{survive} competition, to \emph{grow}, and to \emph{satisfy its customers}. When one of these requirements is missing or poorly handled, the business suffers — exactly like Mama Ngo Bell, whose missing stock records make her lose customers every Saturday.

\begin{definition}[Business need]
A \cle{business need} is something a business \textbf{requires} in order to operate, survive, grow or satisfy its customers. A need may concern materials, money, people, skills, equipment or — crucially for us — \textbf{information and its management}.
\end{definition}

\textbf{The five categories of business needs.} To analyse any business systematically, we group its needs into five families:

\begin{enumerate}
\item \textbf{Operational needs} — everything required to produce the good or deliver the service: raw materials, machines, and above all \emph{stock management} (knowing what is in store, what is running out, what is expiring).
\item \textbf{Administrative needs} — the internal organisation: keeping records, managing staff files, preparing \emph{payroll}, tracking attendance, writing official documents.
\item \textbf{Commercial needs} — everything that links the business to its market: \emph{marketing}, advertising, sales records, and \emph{customer relations} (knowing who buys what, and keeping them loyal).
\item \textbf{Financial needs} — \emph{accounting} (income, expenses, profit), invoicing, tax declarations, and \emph{payments} (cash, cheques, mobile money).
\item \textbf{Security needs} — protecting the business's \emph{property} (premises, stock, cash) and its \emph{data} (customer lists, accounts, passwords) against theft, loss and damage.
\end{enumerate}

\begin{popfigure}
\begin{center}
\begin{tikzpicture}[scale=1.0]
% central circle
\draw[fill=popblue, thick] (0,0) circle (1.35);
\node[white, align=center, text width=2.2cm] at (0,0) {\textbf{Business}\\ \textbf{needs}};
% five branches
\foreach \a/\name/\ex in {90/{Operational}/{Bakery: daily stock of flour},
                          162/{Administrative}/{School: payroll \& staff files},
                          234/{Financial}/{Shop: accounts \& MoMo payments},
                          306/{Security}/{Cybercaf\'e: data backup},
                          18/{Commercial}/{Cooperative: reaching buyers}} {
  \draw[very thick, poporange] (0,0) -- (\a:2.1);
  \node[circle, fill=poporange, draw=popdark, thick, align=center, text width=2.0cm, minimum size=1.7cm, font=\small\bfseries] at (\a:3.05) {\name};
  \node[align=center, text width=3.1cm, font=\footnotesize] at (\a:4.9) {\ex};
}
\end{tikzpicture}
\end{center}
\legende{The five categories of business needs, each illustrated with a Cameroonian example.}
\end{popfigure}

\begin{exemplebox}[Needs of three Cameroonian organisations]
\begin{itemize}
\item A \textbf{bakery in Buea} needs daily sales records to know how many loaves to bake tomorrow (operational + financial need).
\item A \textbf{farmers' cooperative in Bafoussam} needs to reach buyers in Douala and inform them of available quantities and prices (commercial need).
\item A \textbf{secondary school} needs to track fee payments per student and per term, and to send reminders to parents (administrative + financial need).
\end{itemize}
Notice that in each case, the need is fundamentally about \emph{information}: recording it, transmitting it, or analysing it. That is precisely where IT enters.
\end{exemplebox}

\courssub{Problem, Need, Opportunity: Three Words You Must Not Confuse}

GCE examiners regularly test the distinction between these three notions. They form a logical chain:

\begin{definition}[Problem, need and opportunity]
\begin{itemize}
\item A \cle{problem} is an \textbf{existing difficulty}: something that is going wrong or causing a loss right now.
\item A \cle{need} is \textbf{what is required} to solve that problem or to keep the business running well.
\item A \cle{business opportunity} is an \textbf{unmet need in the market} that someone can exploit — a gap that no existing business fills well, and which an entrepreneur can turn into a new product, service or company.
\end{itemize}
\end{definition}

The logic is: \emph{every problem hides a need; every unmet need is a potential opportunity.} Mama Ngo Bell's \emph{problem} is that she runs out of tomatoes before noon. Her \emph{need} is a way to record sales and predict demand. Sandrine's success reveals the \emph{opportunity}: many traders in Mokolo have the same unmet need, so a young graduate could build a business selling a simple stock-tracking service to them.

\begin{methode}[Distinguishing a problem from a need]
When you analyse a business situation, proceed as follows:
\begin{enumerate}
\item Ask ``\emph{What is wrong?}'' — the answer describes the \textbf{problem} (e.g. ``customers leave because tomatoes finish early'').
\item Then ask ``\emph{What is lacking?}'' — the answer states the \textbf{need} (e.g. ``a record of daily sales and a forecast of demand'').
\item Finally ask ``\emph{Who else has this need, and is anyone serving them?}'' — if many people share the need and nobody serves them well, you have found an \textbf{opportunity}.
\end{enumerate}
The key habit is to ask \textbf{``What is lacking?''} and not stop at \textbf{``What is wrong?''}. The first question points to a solution; the second only describes pain.
\end{methode}

\begin{popfigure}
\begin{center}
\begin{tikzpicture}[
  box/.style={rectangle, rounded corners, draw=popdark, very thick, fill=popblueL, align=center, text width=2.6cm, minimum height=1.5cm, font=\small},
  ex/.style={rectangle, draw=popmuted, dashed, align=center, text width=2.6cm, font=\footnotesize, below=0.25cm},
  arr/.style={->, very thick, poporange}]
\node[box, align=center] (p) at (0,0){\textbf{Problem}\\ existing difficulty};
\node[box, align=center] (n) at (4.2,0){\textbf{Need}\\ what is required};
\node[box, align=center] (s) at (8.4,0){\textbf{IT solution}\\ tool that meets the need};
\node[box, align=center] (o) at (12.6,0){\textbf{Opportunity}\\ unmet need in the market};
\draw[arr] (p) -- (n);
\draw[arr] (n) -- (s);
\draw[arr] (s) -- (o);
\node[ex] at (0,-1.6) {Trader runs out of tomatoes before noon};
\node[ex] at (4.2,-1.6) {Record sales and predict demand};
\node[ex] at (8.4,-1.6) {Smartphone stock app + sales log};
\node[ex] at (12.6,-1.6) {Sell the service to all Mokolo traders};
\end{tikzpicture}
\end{center}
\legende{From problem to opportunity: the market trader's stock problem worked through the whole chain.}
\end{popfigure}

\begin{attention}[Common mistake: confusing a business need with the owner's personal want]
A \textbf{business need} serves the \emph{business}: its customers, its operations, its survival. A \textbf{personal want} serves only the owner's comfort or prestige. ``The business needs a record-keeping system so that no sale is forgotten'' is a genuine need. ``The owner wants the latest luxury smartphone to impress friends'' is a personal want — even if the phone is bought with business money. In an examination, always justify a need by its effect on \emph{customers, operations or profit}, never by the owner's desires.
\end{attention}

\courssub{Information: The Raw Material of Good Management}

Why does an ICT course insist so much on business needs? Because behind almost every business need there hides an \emph{information need}. Consider what a manager actually does all day: decide how much stock to order, decide which customers owe money, decide whether the business made a profit this month, decide whether to employ one more worker. \textbf{Every good decision depends on information}, and bad information produces bad decisions as surely as bad flour produces bad bread.

\begin{aretenir}[Information is a business resource]
\cle{Information} is a resource of the business \textbf{just like money, staff and equipment}. A business can own machines and cash and still fail if it does not \emph{know} its stock levels, its debts, its best customers and its real costs. To be useful for decision-making, information must be:
\begin{itemize}
\item \textbf{Accurate} — correct and free of errors (a wrong stock figure leads to a wrong order);
\item \textbf{Timely} — available when needed (yesterday's sales known \emph{this} morning, not next month);
\item \textbf{Relevant} — related to the decision at hand, without useless detail;
\item \textbf{Complete} — nothing important missing.
\end{itemize}
IT (computers, spreadsheets, databases, mobile applications, networks) is the most powerful set of tools ever invented for capturing, storing, processing and communicating such information — which is why spotting business needs is the first skill of the IT entrepreneur.
\end{aretenir}

\begin{exoresolu}[From a real situation to a stated need]
\textbf{Statement.} A private secondary school in Douala keeps fee payments in a paper notebook. At the end of each term, the bursar spends three days checking who has paid, makes frequent errors, and some parents are wrongly accused of not paying. (a) Identify the problem. (b) State the business need. (c) Indicate the category of the need and the type of information the school requires.
\tcblower
\textbf{Solution.}
(a) \emph{Problem} (what is wrong?): fee records kept on paper are slow to consult and full of errors, causing disputes with parents and wasted days of work.

(b) \emph{Need} (what is lacking?): a reliable, quickly searchable record of every payment, showing for each student the amount paid, the balance due and the date of payment.

(c) \emph{Category}: administrative and financial need. \emph{Information required}: accurate (correct amounts), timely (available immediately when a parent asks) and complete (every student, every term). This is exactly the kind of need that a simple database or spreadsheet application can satisfy — and, since many schools share it, it is also a business opportunity for an IT entrepreneur.
\end{exoresolu}

\begin{savaistu}[Mobile Money changed the information game in Cameroon]
Services such as MTN Mobile Money and Orange Money did not only move cash — they created \emph{automatic records}. A trader paid by Mobile Money has, without any extra effort, a dated list of transactions: who paid, how much and when. Many small Cameroonian businesses discovered better financial management simply by reading their mobile money statements. This is a perfect illustration of the chapter's message: when information becomes accurate and timely, decisions improve and businesses grow.
\end{savaistu}

\begin{exercice}[Analysing a business you know]
Choose one small business in your neighbourhood (a call box, a bakery, a hairdressing salon, a ``tournedos'' stand). (a) State what goods or services it provides and who its customers are. (b) List one need in each of the five categories (operational, administrative, commercial, financial, security). (c) For one of these needs, write the corresponding problem as a sentence beginning ``What is wrong is that\ldots'' and the need as a sentence beginning ``What is lacking is\ldots''. (d) Say whether this need, if unmet in many similar businesses, could be a business opportunity, and justify your answer in two sentences.
\end{exercice}

In this section we have built the foundations: a business provides goods or services to customers for money with a goal; a business need is what it requires to operate, survive, grow and satisfy customers; needs fall into five categories; and a problem, a need and an opportunity form a logical chain that the IT entrepreneur learns to read. In the next section, we shall study systematic \emph{techniques} for identifying these needs in a real business — observation, interviews, questionnaires and document analysis.

\coursec{Techniques for Identifying Business Needs}

In the previous section, we explored what a business is, its core functions, and the generic categories of needs it may have. We now move from \emph{what} needs exist to \emph{how} we discover them in a real-world setting. Guessing what a business needs is a recipe for failure; a solution that does not address the actual pain point will be rejected, no matter how technically brilliant it is. This section equips you with a toolkit of systematic techniques to uncover, verify, and prioritise business needs — skills that are examinable and, more importantly, essential for any aspiring ICT entrepreneur or analyst in Cameroon.

\begin{definition}[Needs Identification (Needs Assessment)]
Needs identification is the systematic process of discovering, documenting, and prioritising the gaps between a business's current state and its desired state. It answers the question: \emph{``What problems must be solved, and for whom?''} before any technical solution is proposed.
\end{definition}

\begin{aretenir}[Why Guessing Fails]
\begin{itemize}
    \item \textbf{Wasted resources:} Building a mobile app for a market where traders only use USSD codes.
    \item \textbf{Low adoption:} Staff ignore a new stock system because it does not match their daily workflow.
    \item \textbf{Missed opportunities:} The real bottleneck (e.g., slow cash reconciliation) remains untouched while a flashy website is built.
\end{itemize}
\end{aretenir}

\courssub{The Needs Identification Process: An Overview}

Before diving into each technique, let us visualise the overall workflow. In practice, these techniques are rarely used in isolation; they complement each other.

\begin{popfigure}
\begin{tikzpicture}[node distance=1.8cm, >=stealth, font=\small]
    \tikzset{
        proc/.style={rectangle, rounded corners, draw=popblue, thick, fill=popblueL, text width=2.2cm, align=center, minimum height=1.1cm},
        arrow/.style={->, thick, draw=popdark}
    }
    \node[proc, align=center] (obs){Direct\\Observation};
    \node[proc, right=of obs] (int) {Interviews};
    \node[proc, right=of int, align=center] (surv){Questionnaires\\ \& Surveys};
    \node[proc, below=of int, align=center] (doc){Document\\Analysis};
    \node[proc, right=of doc] (bench) {Benchmarking};
    \node[proc, right=of bench, align=center] (comp){Complaints\\ \& Feedback};
    \node[proc, below=1.2cm of bench, align=center] (prio){Prioritise\\Needs};
    \node[proc, below=1.2cm of prio, align=center] (prob){Write Problem\\Statement};

    \draw[arrow] (obs) -- (int);
    \draw[arrow] (int) -- (surv);
    \draw[arrow] (surv) -- (doc);
    \draw[arrow] (doc) -- (bench);
    \draw[arrow] (bench) -- (comp);
    \draw[arrow] (comp) -- (prio);
    \draw[arrow] (prio) -- (prob);
    % Feedback loops (optional, dashed)
    \draw[arrow, dashed, bend left=20] (prob.west) to node[left, font=\tiny, text=popmuted] {Iterate if needed} (obs.west);
\end{tikzpicture}
\legende{The iterative needs identification process. Start with observation, gather data through multiple sources, then prioritise and articulate the problem.}
\end{popfigure}

\courssub{Direct Observation: ``Go and See'' (Gemba Walk)}

There is no substitute for being physically present where the work happens. In Cameroon, this might mean standing at a \emph{Marché Central} stall, sitting in a \emph{bureau de change} in Douala, or watching a mobile money agent in Buea.

\begin{methode}[Conducting Effective Direct Observation]
\begin{enumerate}
    \item \textbf{Seek permission:} Explain you are a student/analyst studying workflow, not an auditor.
    \item \textbf{Blend in:} Dress appropriately; do not disrupt the flow.
    \item \textbf{Watch for:}
    \begin{itemize}
        \item \textbf{Delays \& queues:} Customers waiting while the seller searches a paper ledger.
        \item \textbf{Workarounds:} A trader using a calculator and a separate notebook because the POS system is too slow.
        \item \textbf{Repeated errors:} Wrong change given, items sold but not recorded.
        \item \textbf{Paper-based tasks:} Carbon-copy receipt books, handwritten stock cards.
        \item \textbf{Physical strain:} Carrying heavy ledgers, walking to a cybercafé to print statements.
    \end{itemize}
    \item \textbf{Take timestamped notes:} ``10:15 -- Agent checks balance on phone, writes on paper, re-enters in Excel later.''
    \item \textbf{Photograph (with consent):} Layout of the shop, forms used, queue length.
\end{enumerate}
\end{methode}

\begin{attention}[Observation Trap]
Do not interview while observing. Observation is for \emph{seeing what people actually do}; interviews are for \emph{understanding why}. Mixing them biases the observation.
\end{attention}

\courssub{Interviews: Structured Conversations}

Observation tells you \emph{what} happens; interviews reveal \emph{why}. You must talk to three groups: \textbf{owners/managers} (strategic view), \textbf{front-line staff} (operational pain), and \textbf{customers} (external perspective).

\begin{methode}[Steps for a Short Business Interview (15--20 minutes)]
\begin{enumerate}
    \item \textbf{Prepare a guide:} Write 5--7 open-ended questions. Example: ``Walk me through what happens when a customer wants to buy on credit.''
    \item \textbf{Start broad, then narrow:} ``How is business generally?'' $\rightarrow$ ``Where do you lose most time?''
    \item \textbf{Use the funnel technique:}
    \begin{itemize}
        \item \textbf{Open:} ``Describe a typical morning.''
        \item \textbf{Probing:} ``You mentioned writing in two books — why both?''
        \item \textbf{Closed (for facts):} ``How many credit sales per day?'' ``Do you use WhatsApp for orders?''
    \end{itemize}
    \item \textbf{Record objectively:} Use a notebook or voice recorder (with permission). Write \emph{verbatim} quotes for key insights.
    \item \textbf{Close with the ``magic wand'' question:} ``If you could change one thing instantly, what would it be?''
    \item \textbf{Thank and summarise:} ``So the main issue is tracking credit sales...''
\end{enumerate}
\end{methode}

\begin{exemplebox}[Interview Snippet: Mobile Money Agent, Yaoundé]
\textbf{Analyst:} ``What is the hardest part of your day?''\\
\textbf{Agent:} ``Reconciling the till. I serve 80 clients. At 6 pm, I count cash, check SMS logs on two phones, and match with the operator's portal. If one transaction is missing, I stay 45 extra minutes.''\\
\textbf{Insight:} The need is \emph{automated daily reconciliation}, not a new customer-facing app.
\end{exemplebox}

\courssub{Questionnaires and Surveys: Scaling Up}

When you need data from many people (e.g., 50 traders in a market), interviews become impractical. A well-designed questionnaire gives you quantifiable evidence.

\begin{methode}[Designing a Mini-Survey]
\begin{enumerate}
    \item \textbf{Define the target population:} ``All clothing retailers in Mokolo Market.''
    \item \textbf{Choose a sample:} At least 10--15 respondents for a class project; use systematic sampling (every 5th stall).
    \item \textbf{Keep it short:} Maximum 10 questions, 5 minutes to complete.
    \item \textbf{Question types:}
    \begin{itemize}
        \item \textbf{Likert scale:} ``How difficult is stock-taking?'' (1 Very easy -- 5 Very hard)
        \item \textbf{Multiple choice:} ``How do you currently record sales?'' (Paper book / Excel / POS / Phone notes / Other)
        \item \textbf{One open-ended:} ``What is your biggest headache?''
    \end{itemize}
    \item \textbf{Pilot test:} Give to 2--3 friends first; fix ambiguous wording.
    \item \textbf{Deploy:} Paper forms (low literacy contexts) or digital tools: \textbf{Google Forms} (free, works offline with Chrome), \textbf{KoboCollect} (designed for field data collection, used by NGOs in Cameroon), \textbf{WhatsApp} (share link in trader groups).
    \item \textbf{Analyse:} Count frequencies, calculate averages, spot patterns.
\end{enumerate}
\end{methode}

\begin{popfigure}
\begin{center}
\begin{tabularx}{\linewidth}{|l|c|c|X|}
\hline
\rowcolor{popblueL}
\textbf{Problem Reported} & \textbf{Frequency (out of 10)} & \textbf{Priority (1=High)} & \textbf{Representative Quote} \\
\hline
Difficulty tracking credit sales & 8 & 1 & ``I forget who owes me; I lose 20k/month.'' \\
\hline
Slow end-of-day reconciliation & 7 & 1 & ``Takes 1 hour; I close late.'' \\
\hline
No digital record of stock & 6 & 2 & ``I guess what to reorder.'' \\
\hline
Customers ask for mobile money receipts & 5 & 2 & ``They want SMS proof.'' \\
\hline
Power cuts corrupt Excel file & 3 & 3 & ``Laptop battery dead; file damaged.'' \\
\hline
Need online presence & 2 & 4 & ``Young customers ask for Instagram.'' \\
\hline
\end{tabularx}
\end{center}
\legende{Extract of a mini-survey of 10 traders in a Cameroonian market. Frequency and priority guide the analyst.}
\end{popfigure}

\courssub{Document Analysis: Mining Existing Records}

Businesses already generate data — you just have to read it. This technique is non-intrusive and reveals \emph{actual} behaviour, not reported behaviour.

\begin{itemize}
    \item \textbf{Receipt books / Sales invoices:} Look for gaps in numbering (missing sales?), frequent corrections, illegible handwriting.
    \item \textbf{Stock cards / Bin cards:} Compare physical count with card balance; large variances indicate theft or poor recording.
    \item \textbf{Account ledgers / Cash books:} Check for arithmetic errors, missing dates, mixing personal and business expenses.
    \item \textbf{Bank / Mobile Money statements:} Identify transaction fees, failed transactions, reconciliation gaps.
    \item \textbf{WhatsApp / SMS logs:} Many Cameroonian SMEs run order management on WhatsApp. Export chat history to analyse order volume, response time, disputes.
\end{itemize}

\begin{exemplebox}[Document Analysis Discovery]
Analysing three months of a boutique's handwritten sales book revealed:
\begin{itemize}
    \item 12\% of entries had no price (``forgot to write'').
    \item 5\% had no date.
    \item The owner could not produce a monthly sales total without spending 2 hours adding manually.
\end{itemize}
\textbf{Need:} A simple point-of-sale that forces mandatory fields and auto-totals.
\end{exemplebox}

\courssub{Benchmarking: Learning from Peers}

``What is the shop next door doing?'' If a competitor uses a \emph{Kentia POS} or \emph{Wave} for payments and serves customers twice as fast, that is a proven need-solution pair.

\begin{methode}[Quick Benchmarking Steps]
\begin{enumerate}
    \item Identify 2--3 similar businesses that appear more efficient.
    \item Visit as a \textbf{mystery shopper}: buy a small item, observe the process.
    \item Note: hardware used, software interface, payment options, customer flow time.
    \item Ask the owner (casually): ``How has the system changed your work?''
    \item Compare with your target business: gaps become needs.
\end{enumerate}
\end{methode}

\courssub{Analysing Complaints and Customer Feedback}

Complaints are \emph{free needs identification data}. Sources:
\begin{itemize}
    \item Verbal complaints at the counter (``You gave me wrong change last time!'').
    \item WhatsApp voice notes from customers.
    \item Google Maps / Facebook reviews (for businesses with a page).
    \item Returned goods / cancelled orders.
\end{itemize}
\textbf{Technique:} Keep a \emph{Complaint Log} for two weeks. Categorise: \emph{Billing error, Stock-out, Slow service, Rude staff, No receipt}. The top category is a high-priority need.

\courssub{Prioritising Needs: From Long List to Short List}

You will uncover 20+ problems. You cannot solve all at once. Use a simple scoring matrix.

\begin{propriete}[Prioritisation Criteria (Weighted Scoring for Upper Sixth)]
For each identified need, score 1--5 on:
\begin{itemize}
    \item \textbf{Urgency:} How soon must it be solved? (Regulatory deadline, imminent loss)
    \item \textbf{Frequency:} How often does the problem occur? (Daily = 5, Monthly = 1)
    \item \textbf{Impact (People affected):} How many staff/customers suffer? (All = 5, One = 1)
    \item \textbf{Potential Gain:} Estimated time/money saved or revenue gained if solved.
\end{itemize}
\textbf{Total Score = Sum of four criteria.} Rank descending. Top 3 = \textbf{Must Have} for the project scope.
\end{propriete}

\begin{exoresolu}[Prioritisation Exercise]
A team identified five needs for a restaurant in Bamenda. Scores (1--5 each):

\begin{center}
\begin{tabular}{|l|cccc|c|}
\hline
\rowcolor{popblueL}
\textbf{Need} & \textbf{Urgency} & \textbf{Freq.} & \textbf{Impact} & \textbf{Gain} & \textbf{Total} \\
\hline
No digital order tracking (waiters shout orders) & 5 & 5 & 5 & 5 & 20 \\
Manual ingredient stock-out & 4 & 4 & 4 & 4 & 16 \\
No card/MoMo payment option & 3 & 3 & 5 & 4 & 15 \\
Website not mobile-friendly & 2 & 2 & 3 & 3 & 10 \\
Owner wants a loyalty app & 1 & 1 & 2 & 2 & 6 \\
\hline
\end{tabular}
\end{center}

\textbf{Solution:} The project scope focuses on \textbf{Need 1} (order tracking) and \textbf{Need 2} (stock alerts). Needs 3--5 are parked for Phase 2.
\end{exoresolu}

\courssub{Writing a Clear Problem Statement}

A problem statement is the \emph{deliverable} of needs identification. It must be understandable by a non-technical stakeholder and precise enough to guide design.

\begin{methode}[The 5W Technique for Problem Statements]
Answer each W explicitly; then synthesise into one paragraph.

\begin{itemize}
    \item \textbf{Who} is affected? (Role, department, number of people)
    \item \textbf{What} is the problem? (Specific symptom, not a solution)
    \item \textbf{Where} does it occur? (Physical location, system, process step)
    \item \textbf{When} does it occur? (Time of day, frequency, trigger event)
    \item \textbf{Why} does it matter? (Consequence: lost revenue, errors, compliance risk, customer dissatisfaction)
\end{itemize}
\end{methode}

\begin{exemplebox}[From 5W to Problem Statement]
\textbf{Scenario:} A wholesale distributor in Douala.
\begin{itemize}
    \item \textbf{Who:} 3 sales clerks and the accountant.
    \item \textbf{What:} Credit sales are recorded in a paper ledger; no real-time visibility of outstanding debts.
    \item \textbf{Where:} At the sales counter and in the back office.
    \item \textbf{When:} Daily; accountant reconciles only on Fridays.
    \item \textbf{Why:} 15\% of credit sales are never collected (est. 500,000 FCFA/month loss); disputes with customers over amounts owed.
\end{itemize}
\textbf{Problem Statement:}
``The three sales clerks at the Douala wholesale counter record credit sales in a paper ledger, giving the accountant no real-time visibility of outstanding debts. Because reconciliation occurs only weekly, an estimated 15\% of credit value (≈500,000 FCFA/month) is never collected, and frequent disputes arise with customers over amounts owed.''
\end{exemplebox}

\begin{attention}[Common Mistake: Solutionising Too Early]
\textbf{Wrong:} ``The business needs a mobile app to track credit sales.'' \\
\textbf{Right:} ``The business needs a way to record and monitor credit sales in real time to reduce uncollected debts.'' \\
The first prescribes a technology (mobile app); the second states the \emph{need}. The solution could be a USSD menu, a simple web form, or an offline-first Android app — let feasibility decide.
\end{attention}

\begin{aretenir}[Exam Tip: GCE Scenario Questions]
\begin{itemize}
    \item GCE questions often present a short case study (e.g., ``Mr. Nfor runs a pharmacy...'').
    \item \textbf{Task:} ``State the business need.'' or ``Identify two problems with the current system.''
    \item \textbf{Technique:} \textbf{Quote evidence from the text.} Do not invent.
    \item \textbf{Example answer:} ``The business need is to reduce dispensing errors. Evidence: the scenario states `the pharmacist manually writes labels and has made three dosage errors this month'.''
    \item Always link the need to a \emph{consequence} (loss, risk, inefficiency).
\end{itemize}
\end{aretenir}

\courssub{Summary Checklist for Your Project}

Before moving to Section 3 (Matching Needs to IT Solutions), ensure you have:
\begin{enumerate}
    \item Spent at least 1 hour observing the business (photos + notes).
    \item Conducted 2--3 interviews (owner + 1 staff + 1 customer if possible).
    \item Distributed and collected 10+ questionnaires (paper or Google Forms/KoboCollect).
    \item Analysed 2--3 key documents (ledger, stock card, WhatsApp logs).
    \item Benchmarked 1 competitor.
    \item Compiled a prioritised list (Weighted Scoring scores).
    \item Written \textbf{one} polished problem statement using the 5W method.
\end{enumerate}

This evidence portfolio is what examiners expect to see in your project documentation and what real investors or partners will ask for. In the next section, we will explore how to map these validated needs to appropriate IT solutions — from simple spreadsheets to custom software — while respecting the Cameroonian context of connectivity, cost, and digital literacy.

\coursec{Matching Business Needs to IT Solutions}

In the previous section we learned how to \emph{identify} business needs: observing the workplace, interviewing workers and customers, analysing documents and listing problems. But identifying a need is only half of the work. A need written on paper does not solve anything by itself. The entrepreneur must now answer a second question: \textbf{which IT tool can satisfy this need?} This section teaches you how to build that bridge between a clearly stated need and an appropriate IT solution.

\courssub{The Guiding Principle: Technology Is a Means, Not an End}

Imagine a tailor in Buea who keeps losing customers' measurement books. A friend advises him to buy a powerful computer with the latest design software. The computer arrives, but the tailor does not know how to use design software, has no electricity stability for it, and his real problem --- lost notebooks --- remains unsolved. A simple smartphone application for storing customer measurements would have solved the problem for a fraction of the cost.

This story illustrates the golden rule of this chapter: \textbf{we never start from the technology; we always start from the need.} Technology is a \emph{means} to satisfy a need, never an \emph{end} in itself.

\begin{definition}[IT solution (business context)]
An \cle{IT solution} is a combination of hardware, software and services (computers, phones, applications, networks, training) put in place to satisfy one or more clearly identified business needs. An IT solution is said to be \cle{relevant} only if it actually removes or reduces the problem expressed in the need.
\end{definition}

\begin{attention}[Choosing the most sophisticated instead of the most appropriate]
A very common mistake is to buy the most expensive or most fashionable technology: a big server where a spreadsheet would do, a custom-built application where WhatsApp Business is enough. Sophisticated technology brings hidden costs: maintenance, training, electricity, internet subscription. The \emph{best} solution is the \textbf{simplest one that fully satisfies the need}, not the most impressive one.
\end{attention}

\begin{propriete}[Criteria of a good IT solution]
A good IT solution for a small business must be:
\begin{itemize}
\item \cle{Relevant}: it directly addresses the identified need;
\item \cle{Affordable}: its total cost (purchase, subscription, maintenance, electricity) fits the budget of the business;
\item \cle{Usable}: the owner and workers can learn to use it quickly, with their existing skills;
\item \cle{Sustainable}: it can be maintained, repaired and updated locally over several years.
\end{itemize}
(Justifying a choice against these four criteria is a classic GCE examination question.)
\end{propriete}

\courssub{The Need-to-Solution Matching Procedure}

Matching a need to a solution is not guesswork; it follows a clear procedure.

\begin{methode}[Matching a need to an IT solution]
\begin{enumerate}
\item \textbf{State the need} precisely, in one sentence (e.g. ``the shopkeeper cannot know his stock level without counting by hand'').
\item \textbf{List possible IT tools} that could address this need (spreadsheet, stock software, mobile app).
\item \textbf{Compare} the candidates using the four criteria: relevance, affordability, usability, sustainability.
\item \textbf{Choose} the simplest tool that satisfies the need within the budget.
\item \textbf{Justify} the choice in writing, linking each feature of the tool to the stated need.
\end{enumerate}
\end{methode}

\begin{popfigure}
\begin{center}
\begin{tikzpicture}[font=\small]
% Left column: needs
\node[fill=popblueL, draw=popblue, rounded corners, text width=3.6cm, align=center, minimum height=0.9cm] (n1) at (0,4.5) {Lost records and difficult calculations};
\node[fill=popblueL, draw=popblue, rounded corners, text width=3.6cm, align=center, minimum height=0.9cm] (n2) at (0,3.0) {Slow cash collection and change problems};
\node[fill=popblueL, draw=popblue, rounded corners, text width=3.6cm, align=center, minimum height=0.9cm] (n3) at (0,1.5) {Customers do not know the products};
\node[fill=popblueL, draw=popblue, rounded corners, text width=3.6cm, align=center, minimum height=0.9cm] (n4) at (0,0.0) {Theft and insecurity at the shop};
% Right column: solutions
\node[fill=popgreenL, draw=popgreen, rounded corners, text width=3.6cm, align=center, minimum height=0.9cm] (s1) at (8,4.5) {Spreadsheet or database for records};
\node[fill=popgreenL, draw=popgreen, rounded corners, text width=3.6cm, align=center, minimum height=0.9cm] (s2) at (8,3.0) {Mobile money payments (MoMo, OM)};
\node[fill=popgreenL, draw=popgreen, rounded corners, text width=3.6cm, align=center, minimum height=0.9cm] (s3) at (8,1.5) {Facebook page and WhatsApp catalogue};
\node[fill=popgreenL, draw=popgreen, rounded corners, text width=3.6cm, align=center, minimum height=0.9cm] (s4) at (8,0.0) {CCTV, passwords and backups};
% Arrows
\draw[->, thick, poporange] (n1) -- (s1);
\draw[->, thick, poporange] (n2) -- (s2);
\draw[->, thick, poporange] (n3) -- (s3);
\draw[->, thick, poporange] (n4) -- (s4);
% Titles
\node[font=\bfseries\small, popdark] at (0,5.4) {Business needs};
\node[font=\bfseries\small, popdark] at (8,5.4) {IT solutions};
\end{tikzpicture}
\end{center}
\legende{Matching diagram: each identified need is linked to an appropriate IT solution.}
\end{popfigure}

\courssub{Matching the Main Categories of Needs to IT Tools}

\textbf{Record keeping and data management.}\par
When a business loses documents, forgets customer debts or cannot find past transactions, the need is \emph{organised data storage}. For simple lists and calculations, a \cle{spreadsheet} such as Microsoft Excel is enough: one sheet for sales, one for expenses, with automatic totals. When data become large and interrelated (customers, products, orders), a \cle{database} such as Microsoft Access, or dedicated \cle{stock management software}, is more appropriate because it avoids duplication and allows fast searches.

\textbf{Accounting and financial needs.}\par
For daily bookkeeping, profit calculation and tax declarations, \cle{accounting software} records every transaction and produces reports automatically. For payments and collections, \cle{mobile money} services (MTN Mobile Money, Orange Money) allow a trader to receive money instantly without handling cash, reducing theft and change problems.

\textbf{Communication needs.}\par
To communicate with suppliers and customers: \cle{email} for formal documents and quotations, \cle{WhatsApp Business} for quick answers, price lists and product photos, \cle{bulk SMS} to announce promotions to many customers at once, and a \cle{website} for a permanent public presence.

\textbf{Marketing and sales needs.}\par
To reach more customers: \cle{social media pages} (Facebook, Instagram) to display products, \cle{e-commerce platforms} (Jumia, Glotelho in Cameroon) to sell online with delivery, and \cle{online catalogues} shared through WhatsApp.

\textbf{Security needs.}\par
To protect money, data and premises: strong \cle{passwords} on devices and accounts, regular \cle{backups} of records (USB drive or cloud), an \cle{antivirus}, \cle{CCTV cameras} connected to a computer, and \cle{access control} so that only authorised staff open sensitive files.

\textbf{Education and training needs.}\par
Schools and training centres use \cle{e-learning platforms} and digital content to teach beyond the classroom, and \cle{school management software} to handle marks, report cards and fees.

\textbf{Agriculture and transport needs.}\par
Farmers receive \cle{market prices by SMS or apps} before selling, avoiding cheating by middlemen. Transporters use \cle{GPS tracking} to follow vehicles, and \cle{ride-hailing and logistics apps} connect drivers to customers and goods.

\begin{aretenir}[The mobile phone: the most realistic platform]
For most small Cameroonian businesses, the \cle{smartphone} is the most realistic IT platform: it is already owned, cheap to run, works with mobile money, WhatsApp, Facebook, SMS and many business apps, and does not depend on stable electricity. Always consider a phone-based solution before proposing a computer-based one.
\end{aretenir}

\begin{popfigure}
\begin{center}
\begin{tikzpicture}[font=\small, grow cyclic, mindmap, every node/.style={concept, circular drop shadow},
root concept/.append style={concept color=popblue, text=white, font=\bfseries\small},
level 1/.append style={level distance=3.4cm, sibling angle=72},
level 1 concept/.append style={concept color=poporange, text width=2.4cm, align=center, font=\small},
level 2/.append style={level distance=2.4cm, sibling angle=45},
level 2 concept/.append style={concept color=popgreenL, text width=2.0cm, align=center, font=\footnotesize}]
\node[root concept, text width=3.0cm, align=center] {IT solutions for a small Cameroonian shop}
child { node {Sales} child { node {Excel sales sheet} } child { node {Daily totals} } }
child { node {Stock} child { node {Stock software} } child { node {Low-stock alerts} } }
child { node {Payments} child { node {MTN MoMo} } child { node {Orange Money} } }
child { node {Marketing} child { node {Facebook page} } child { node {WhatsApp catalogue} } }
child { node {Security} child { node {Passwords \& backups} } child { node {CCTV} } };
\end{tikzpicture}
\end{center}
\legende{Mind map of IT solutions for a small Cameroonian shop, organised by business function.}
\end{popfigure}

\courssub{Case Analysis Table}

The following table summarises the matching for typical Cameroonian situations. Learn to produce such a table yourself in the examination.

\begin{center}
\begin{tabularx}{\linewidth}{|l|l|X|X|}
\hline
\rowcolor{popblueL} \textbf{Need} & \textbf{IT tool} & \textbf{Expected benefit} & \textbf{Cameroonian example} \\
\hline
Record keeping & Excel spreadsheet & Fast, error-free calculations & A Douala provisions store totals daily sales automatically \\
\hline
Stock control & Stock management software & No more stock-outs or over-buying & A pharmacy in Yaound\'e is alerted when a medicine runs low \\
\hline
Cash collection & MTN MoMo / Orange Money & Instant, safe payments & A Buea snack owner receives payments without handling cash \\
\hline
Customer communication & WhatsApp Business & Quick replies, product photos & A tailor in Bamenda sends design photos to clients \\
\hline
Marketing & Facebook / Instagram page & Reaches thousands of customers cheaply & A caterer in Limbe advertises cakes online \\
\hline
Online sales & Jumia, Glotelho & Sells beyond the local quarter & An electronics seller delivers nationwide \\
\hline
Security & Passwords, backups, CCTV & Protects money, data and goods & A cybercaf\'e keeps CCTV linked to its main computer \\
\hline
Farm information & Price SMS / apps & Better selling prices & A cocoa farmer checks prices before selling \\
\hline
Transport & GPS tracking, logistics apps & Follows vehicles, finds clients & A transport agency tracks its buses on the Douala--Yaound\'e road \\
\hline
\end{tabularx}
\end{center}

\begin{exoresolu}[Choosing an IT solution for a bookshop]
\textbf{Statement.} Madam Ekane runs a bookshop in Yaound\'e. She complains: ``I never know exactly how many copies of each textbook remain; I discover shortages only when a customer asks.'' Her budget is small and she owns a smartphone and an old laptop. Propose an IT solution and justify it using the four criteria.
\tcblower
\textbf{Solution.} \emph{Step 1 --- State the need:} the shopkeeper needs to know stock levels in real time without counting by hand. \emph{Step 2 --- Possible tools:} paper register (rejected, already failing), Excel spreadsheet, dedicated stock software. \emph{Step 3 --- Compare:} dedicated stock software is powerful but costly and requires training; Excel is cheap, already available on her laptop, and sufficient for a single shop. \emph{Step 4 --- Choose:} an Excel workbook with one row per title, columns for quantity in, quantity sold and remaining stock computed by formula. \emph{Step 5 --- Justify:} \textbf{relevant} (each sale updates the remaining stock, so shortages are visible immediately); \textbf{affordable} (uses the existing laptop, no subscription); \textbf{usable} (basic Excel skills suffice); \textbf{sustainable} (the file can be backed up on her phone or a USB drive and improved later). If the business grows, she can migrate to stock management software.
\end{exoresolu}

\begin{exercice}[Match and justify]
A maize farmer in the North-West Region sells his harvest to middlemen at very low prices because he never knows the real market prices in the towns. (a) State his need in one sentence. (b) Propose two possible IT tools. (c) Choose the most appropriate one and justify your choice with the four criteria of a good IT solution.
\end{exercice}

\begin{corrige}[Exercise 1]
(a) Need: the farmer needs access to current market prices before selling his maize. (b) Possible tools: a price-information service by SMS or a smartphone application; alternatively a computer with internet access at the cooperative. (c) Choice: the SMS/app service on his mobile phone. Justification: \textbf{relevant} --- it delivers prices directly before negotiation; \textbf{affordable} --- a basic phone and cheap SMS subscription; \textbf{usable} --- no advanced skills required; \textbf{sustainable} --- works with the mobile networks already covering rural areas, unlike a computer which needs electricity and internet.
\end{corrige}

\coursec{Evaluating Feasibility and Spotting IT Business Opportunities}

In the previous section we learned how to match an identified business need to a suitable IT solution. But proposing a solution on paper is not enough. A shop owner in Bamenda may dream of an online store, yet have no reliable electricity; a school in Ngaoundér\'e may want a computerised report-card system, yet have no trained staff. Before any money is spent, one question must be answered honestly: \emph{can this solution actually work here, for this business, at an acceptable cost?} That is the work of a \cle{feasibility study}. This final section teaches you how to evaluate feasibility along four dimensions, how to turn an unsolved need into a genuine IT business opportunity, and how to write a short feasibility report — a skill directly tested at the GCE Advanced Level.

\courssub{Why Feasibility Matters}

\begin{experience}[A story repeated across Cameroon]
A young entrepreneur in Bafoussam buys ten computers to open a modern secretarial bureau. Within six months the business closes. Why? The electricity cuts damaged two power supplies, the internet subscription cost more than expected, customers were fewer than hoped, and no budget had been kept for maintenance. None of these problems was technical in the strict sense: the computers worked. The project failed because \textbf{feasibility was never studied}.
\end{experience}

Many IT projects in Cameroon fail not because the technology is bad, but because of \textbf{cost overruns}, \textbf{electricity cuts}, \textbf{poor connectivity}, or \textbf{lack of digital skills}. A feasibility study exists precisely to detect these obstacles \emph{before} investing.

\begin{definition}[Feasibility study]
A \cle{feasibility study} is a systematic evaluation, carried out \emph{before} implementation, of whether a proposed project or IT solution can realistically succeed in its context. It examines whether the solution is technically possible, economically worthwhile, usable by the people concerned, and acceptable legally and socially.
\end{definition}

\courssub{The Four Dimensions of Feasibility}

\textbf{Technical feasibility.}\par
This asks: \emph{do we have, or can we obtain, the physical and digital means?} Key questions include:
\begin{itemize}
\item Is the required \textbf{hardware} (computers, printers, POS terminals, smartphones) available on the local market, and can spare parts be found?
\item Is the \textbf{software} available, licensed and compatible with the hardware?
\item Is \textbf{electricity} reliable? If not, can the business afford a generator, an inverter, or a solar panel kit with batteries?
\item Is there adequate \textbf{internet coverage} (CAMTEL fibre, MTN or Orange 3G/4G) at the business location, and at what speed?
\end{itemize}

\textbf{Economic feasibility.}\par
This asks: \emph{do the benefits outweigh the costs?} Costs include the purchase of equipment, installation, monthly subscriptions (internet data, software licences), electricity or fuel for a generator, and maintenance. Benefits include time saved, errors reduced, new customers attracted, and extra revenue. A simple \cle{cost-benefit reasoning} compares the estimated monthly cost of the solution with the estimated monthly gain it brings.

\textbf{Human feasibility.}\par
This asks: \emph{can the people actually use the tool?} One must assess the digital skills of the owner and staff, the need (and cost) of \textbf{training}, and the \textbf{ease of use} of the chosen tool. A sophisticated stock-management package is useless if the only employee has never used a spreadsheet; a simple mobile application may be far more feasible.

\textbf{Legal and social feasibility.}\par
This asks: \emph{is the solution lawful and accepted?} In Cameroon, businesses handling personal data must respect data protection principles (ANTIC is the national agency in charge of ICT regulation and electronic security). Socially, one must consider \textbf{customer acceptance}: will customers trust digital records and mobile money payments, or do they still prefer cash and paper receipts? A solution rejected by customers is not feasible, however brilliant it is.

\begin{methode}[The four-dimension feasibility checklist]
To evaluate any proposed IT solution, proceed dimension by dimension:
\begin{enumerate}
\item \textbf{Technical:} list the hardware, software, power and connectivity required; verify local availability and reliability.
\item \textbf{Economic:} estimate all initial and \emph{recurrent} costs; estimate expected gains; compare.
\item \textbf{Human:} assess current skills of users; plan training if needed; prefer simple tools.
\item \textbf{Legal/social:} check data protection obligations and customer acceptance.
\item \textbf{Conclude:} for each dimension, answer Yes / Partly / No. Any ``No'' must be resolved, or the project redesigned, before investing.
\end{enumerate}
\end{methode}

\begin{popfigure}
\begin{center}
\begin{tikzpicture}[scale=0.9]
% Header row
\fill[popblueL] (0,4) rectangle (10.4,5);
\node[anchor=west, font=\bfseries] at (0.1,4.5) {Feasibility dimension};
\node[font=\bfseries] at (6.9,4.5) {Yes};
\node[font=\bfseries] at (8.3,4.5) {Partly};
\node[font=\bfseries] at (9.7,4.5) {No};
% Row labels
\node[anchor=west] at (0.1,3.5) {\textbf{Technical:} hardware, software, power, internet};
\node[anchor=west] at (0.1,2.5) {\textbf{Economic:} costs vs expected benefits};
\node[anchor=west] at (0.1,1.5) {\textbf{Human:} skills, training, ease of use};
\node[anchor=west] at (0.1,0.5) {\textbf{Legal/social:} data protection, acceptance};
% Grid lines
\draw[popdark] (0,0) rectangle (10.4,5);
\draw[popdark] (0,4) -- (10.4,4);
\draw[popdark] (6.2,0) -- (6.2,5);
\draw[popdark] (7.6,0) -- (7.6,5);
\draw[popdark] (9.0,0) -- (9.0,5);
\foreach \y in {1,2,3}{\draw[popline] (0,\y) -- (10.4,\y);}
% Checkboxes (empty)
\foreach \y in {3.5,2.5,1.5,0.5}{
  \draw[popdark] (6.6,\y-0.18) rectangle (7.2,\y+0.18);
  \draw[popdark] (8.0,\y-0.18) rectangle (8.6,\y+0.18);
  \draw[popdark] (9.4,\y-0.18) rectangle (10.0,\y+0.18);
}
% Sample ticks: technical Partly, economic Yes, human Partly, legal Yes
\node[popgreen, font=\bfseries] at (8.3,3.5) {\checkmark};
\node[popgreen, font=\bfseries] at (6.9,2.5) {\checkmark};
\node[popgreen, font=\bfseries] at (8.3,1.5) {\checkmark};
\node[popgreen, font=\bfseries] at (6.9,0.5) {\checkmark};
\end{tikzpicture}
\end{center}
\legende{Feasibility checklist grid: each dimension is ticked Yes, Partly or No (sample answers shown).}
\end{popfigure}

\begin{attention}[Common mistake: ignoring recurrent costs]
When judging whether an IT solution is affordable, students (and real entrepreneurs!) often count only the \emph{purchase price} of the equipment. This is a trap. The true cost includes \textbf{recurrent costs}: monthly internet data bundles, electricity or generator fuel, software subscriptions, repairs and maintenance, and staff training. A computer costing 150\,000 FCFA may cost another 20\,000 FCFA per month to keep running. Always compare \emph{total monthly cost} with \emph{monthly benefit}.
\end{attention}

\begin{exoresolu}[Feasibility of a POS system for a provision store in Douala]
\textbf{Statement.} Mrs~Ekane owns a provision store in Bonabéri, Douala. She wants to replace her paper sales notebook with a point-of-sale (POS) terminal costing 120\,000 FCFA, plus software subscription and internet data of 15\,000 FCFA per month. She estimates the system will save her about 25\,000 FCFA per month by reducing stock losses and calculation errors. Electricity cuts are frequent in her quarter, but she already owns a small generator. Assess the feasibility and advise her.
\tcblower
\textbf{Solution.}
\begin{itemize}
\item \textbf{Technical:} POS terminals are sold in Douala; she has a generator for power cuts; 4G coverage is good in Bonabéri. \emph{Feasible (Yes).}
\item \textbf{Economic:} monthly cost $\approx 15\,000$ FCFA (plus maintenance allowance, say 5\,000 FCFA) $= 20\,000$ FCFA; monthly gain $\approx 25\,000$ FCFA. Net gain $\approx 5\,000$ FCFA/month, so the 120\,000 FCFA investment is recovered in about $120000/5000 = 24$ months. \emph{Feasible (Yes), but the payback period is long.}
\item \textbf{Human:} Mrs~Ekane uses a smartphone daily; a short training session on the POS app should suffice. \emph{Feasible (Yes).}
\item \textbf{Legal/social:} sales records are her own business data; customers receive printed receipts, which increases trust. \emph{Feasible (Yes).}
\end{itemize}
\textbf{Advice:} the project is feasible on all four dimensions, but she should budget for maintenance and accept a payback period of about two years.
\end{exoresolu}

\begin{popfigure}
\begin{center}
\begin{tikzpicture}
\begin{axis}[
  ybar, bar width=0.9cm,
  width=10cm, height=6cm,
  ymin=0, ymax=30,
  ylabel={Amount (thousand FCFA / month)},
  symbolic x coords={POS terminal, Website, Stock app},
  xtick=data,
  legend style={at={(0.5,1.02)}, anchor=south, legend columns=2},
  nodes near coords, every node near coord/.append style={font=\small},
]
\addplot[fill=poporange] coordinates {(POS terminal,20) (Website,12) (Stock app,8)};
\addplot[fill=popgreen] coordinates {(POS terminal,25) (Website,18) (Stock app,6)};
\legend{Estimated monthly cost, Estimated monthly gain}
\end{axis}
\end{tikzpicture}
\end{center}
\legende{Illustrative comparison of estimated monthly cost versus estimated monthly gain for three proposed IT solutions in a small shop (fictitious figures).}
\end{popfigure}

\courssub{From Need to Opportunity: Becoming an IT Entrepreneur}

Here is the entrepreneurial insight: \emph{every unsolved business need you can identify is a potential business for you}. When you notice that shopkeepers in your quarter struggle with stock records, that farmers lack price information, or that students cannot find evening tutors, you have spotted a gap that an IT solution could fill — and you could be the one providing it.

\begin{definition}[IT-based business opportunity]
An \cle{IT-based business opportunity} is an identified, unsolved need of people or businesses that can be met profitably by providing an information technology product or service. It exists when (a) the need is real and widespread, (b) an IT solution is feasible, and (c) customers are willing and able to pay for it.
\end{definition}

Cameroonian examples abound:
\begin{itemize}
\item \textbf{Mobile money agents and MoMo kiosks:} the need for easy money transfer and payment in areas without banks created thousands of MTN MoMo and Orange Money agent businesses.
\item \textbf{Cybercafés and digital service bureaus:} the need for internet access, printing, typing and online administrative procedures (tax declarations, exam registration) sustains these ventures in every town.
\item \textbf{Agro-tech and agri-info services:} farmers need weather forecasts, market prices and veterinary advice; SMS and app-based services answer this need.
\item \textbf{Digital health services:} appointment booking, health information and teleconsultation answer the need for access to care.
\item \textbf{Online tutoring and e-learning:} students preparing for the GCE need affordable revision support — a need answered by tutoring platforms and educational apps.
\end{itemize}

\begin{savaistu}[Did you know?]
The mobile money agent network is one of the largest IT-driven job creators for young Cameroonians. A single kiosk requires only a smartphone or basic phone, a float of cash and e-money, and a good location — a perfect illustration of a \emph{feasible} IT business: low technical barrier, clear economic benefit, minimal training, and strong social acceptance because customers already trust the MTN and Orange brands.
\end{savaistu}

\courssub{Writing a Short Needs Analysis and Feasibility Report}

At the GCE, and in real life, you must be able to present your analysis in a clear, structured document. A short needs analysis and feasibility report contains the following sections:
\begin{enumerate}
\item \textbf{Context:} brief description of the business or community (who, where, what activity).
\item \textbf{Problem statement:} the difficulty observed, stated in one or two precise sentences.
\item \textbf{Identified needs:} the list of needs deduced from the problem.
\item \textbf{Proposed IT solution:} the tool or service chosen, with a short description.
\item \textbf{Feasibility:} the four-dimension assessment (technical, economic, human, legal/social).
\item \textbf{Expected benefits:} measurable gains (time, money, quality, customers) and a conclusion.
\end{enumerate}

\begin{exercice}[Feasibility report]
A cooperative of cassava farmers near Ebolowa wants to know the daily market prices in Yaoundé before sending their produce, and wants to receive mobile money payments from buyers. Write a short needs analysis and feasibility report (six sections) proposing an appropriate IT solution.
\end{exercice}

\begin{corrige}[Exercise 1]
\textbf{Context:} a cassava farmers' cooperative near Ebolowa, selling in Yaoundé markets. \textbf{Problem:} farmers sell blindly, without knowing current prices, and cash handling is risky. \textbf{Needs:} (i) daily price information; (ii) secure payment. \textbf{Proposed solution:} a WhatsApp/SMS price-information group managed by the cooperative secretary, plus MTN MoMo/Orange Money merchant accounts. \textbf{Feasibility:} \emph{technical} — 3G/4G coverage exists, members own phones (Yes); \emph{economic} — cost limited to data bundles of a few thousand FCFA per month, against better sale prices (Yes); \emph{human} — members already use phones; brief training on MoMo needed (Partly); \emph{legal/social} — mobile money is widely trusted (Yes). \textbf{Expected benefits:} better negotiated prices, reduced cash theft risk, faster payments. Conclusion: the project is feasible and recommended.
\end{corrige}

\courssub{GCE Focus: Answering Scenario Questions (Exam Extra)}

Scenario questions at the GCE Advanced Level typically describe a business and ask you to analyse it. Follow this four-step plan:
\begin{enumerate}
\item \textbf{Identify the need:} quote or rephrase the exact difficulty stated in the scenario.
\item \textbf{Propose the IT solution:} name a concrete, appropriate tool (not just ``a computer'').
\item \textbf{Justify the choice:} explain, using feasibility arguments, why this solution fits \emph{this} business.
\item \textbf{State one limitation:} show critical thinking — cost, electricity, connectivity, skills or acceptance.
\end{enumerate}

\begin{attention}[Exam tip: always link back to the scenario]
In justification questions, examiners award marks only when the proposed IT solution is explicitly linked to the \emph{specific need stated in the scenario}. Writing ``a computer is fast and efficient'' earns little; writing ``a stock-management application will alert Mrs~X when flour runs low, which solves her stated problem of frequent stock-outs'' earns full marks. Need $\rightarrow$ solution $\rightarrow$ justification $\rightarrow$ limitation: never break the chain.
\end{attention}

\begin{aretenir}[Key points]
\begin{itemize}
\item A \cle{feasibility study} evaluates, before investment, whether an IT solution can succeed: technically, economically, humanly, and legally/socially.
\item Always include \textbf{recurrent costs} (data, electricity, maintenance, training) in the economic analysis.
\item An unsolved, widespread need + a feasible IT solution + customers willing to pay $=$ an \cle{IT-based business opportunity}.
\item A feasibility report follows six sections: context, problem, needs, solution, feasibility, expected benefits.
\item At the GCE: identify the need, propose the solution, justify by linking to the scenario, state one limitation.
\end{itemize}
\end{aretenir}

\newpage

\coursec{Integration activity}

\coursec{Starting a Business: Identifying IT Needs}

\courssub{Aims of the Activity}

This integration activity, titled \textbf{IT Needs Detective in Our Community}, mobilises all the competences developed in this chapter in one realistic field mission. By the end of the activity, you should be able to:

\begin{itemize}
\item \textbf{Observe} a real business and identify operational problems that may be solvable with IT (Lesson 51);
\item \textbf{Interview} a business owner using a structured guide and extract relevant information;
\item \textbf{Classify} identified needs by category (record keeping, communication, security, marketing, payments, stock management);
\item \textbf{Match} each priority need to an appropriate, affordable IT solution;
\item \textbf{Evaluate} the feasibility of the main proposal along four dimensions: technical, economic, operational and schedule;
\item \textbf{Write} a one-page needs analysis report and \textbf{defend} it in a 5-minute oral presentation.
\end{itemize}

\begin{attention}[Before going to the field]
Always request permission from the business owner before observing or interviewing. Introduce yourself as a student, explain the purpose of the exercise, be polite and discreet, and never promise that a solution will actually be installed. You are \emph{analysing} needs, not selling products.
\end{attention}

\begin{methode}[The 5-Step Needs Analysis Process]
Follow this sequence for every business you study:
\begin{enumerate}
\item \textbf{Observe} silently: watch workflows, note tools, spot bottlenecks.
\item \textbf{Interview} with a prepared guide: mix open and closed questions.
\item \textbf{Classify} problems into standard need categories.
\item \textbf{Match} each category to a realistic IT solution (hardware, software, service).
\item \textbf{Check feasibility} on four dimensions: Technical, Economic, Operational, Schedule (TEOS).
\end{enumerate}
\end{methode}

\courssub{Situation}

\textbf{Mama Caro's Provision Store, Mile 4 -- Bamenda.}\par
\medskip
Mama Caro has run a provision store near your school for nine years. She sells food items (rice, groundnut oil, maggi cubes, soap, sachet tomatoes), phone accessories and prepaid credit transfer for MTN and Orange. The shop opens from 6:30 to 20:30, six days a week, and employs one shop assistant, Blessing. Business is good in the morning and at closing time from school, but Mama Caro is worried. During a first informal chat, she tells your team:

\begin{quote}
\emph{``My children, I work hard but at the end of the month I don't know where the money goes. Everything is written in this exercise book, but Blessing sometimes forgets to record credit sales, and customers deny their debts. Last month I discovered that three cartons of oil had finished without me noticing --- I lost sales for two days. The credit transfer money is mixed with the shop money and I can't separate them. Thieves broke into the shop opposite last term, so I fear for my cash at night. Other shops now paste their prices on WhatsApp status and customers go there. I have a smartphone, electricity most days, and a small generator, but no computer. I can spend at most 150\,000 FCFA this year on improvements.''}
\end{quote}

Your team's 20-minute observation confirms and completes her statements:

\begin{center}
\begin{tabularx}{\linewidth}{|l|X|}\hline
\rowcolor{popblueL} \textbf{Observation} & \textbf{Detail noted by the team} \\ \hline
Record keeping & One paper exercise book; entries in pencil; several pages torn; credit sales written on loose papers. \\ \hline
Stock control & No stock list; Mama Caro checks shelves visually; no alert when items run low. \\ \hline
Payments & Cash only; MTN MoMo and Orange Money transfers recorded in the same book as shop sales. \\ \hline
Security & Cash box under the counter; no camera; no backup of records (book could be stolen or burnt). \\ \hline
Marketing & No price display outside; no customer contact list; competitors advertise on WhatsApp. \\ \hline
Equipment & One Android smartphone (Mama Caro), no computer, no printer; electricity with occasional cuts; small generator available. \\ \hline
\end{tabularx}
\end{center}

The class is divided into teams of four. Your team must produce a complete \textbf{needs analysis report} and present it in 5 minutes. The class will then vote for the most realistic and impactful proposal.

\begin{savaistu}[Mobile Money in Cameroon]
Cameroon has one of the highest mobile money adoption rates in Central Africa. MTN MoMo and Orange Money together serve over 10 million active wallets. For small retailers, becoming a MoMo/OM agent brings commission income but requires strict float separation --- a classic IT need that a simple spreadsheet or agent app can solve.
\end{savaistu}

\courssub{Tasks}

\begin{enumerate}
\item \textbf{Observation and problem listing.} From the observation table and Mama Caro's statement, list at least six distinct problems affecting the business. For each problem, state its visible consequence (loss of money, loss of sales, wasted time, insecurity, etc.).
\item \textbf{Interview guide.} Write the six questions your team would ask Mama Caro to complete the observation. At least two questions must be open questions, and each question must target information you cannot get by simple observation (e.g.\ frequency of stock-outs, size of credit debts, willingness to learn a new tool).
\item \textbf{Classification of needs.} Group the problems from question 1 into need categories (record keeping and accounting, stock management, payments and separation of funds, security, marketing and customer relations). Present your classification in a table with three columns: \emph{Need category -- Problem observed -- Evidence from the field}.
\item \textbf{Matching needs to IT solutions.} For each need category, propose one IT solution that is realistic for this shop (hardware, software or service). Justify each match in one sentence, taking into account the available equipment (smartphone, no computer, generator) and the budget of 150\,000 FCFA.
\item \textbf{Priority and feasibility.} Choose the \emph{one} proposal you consider the priority. Fill the four-dimension feasibility checklist: \textbf{technical} (does the technology exist and can it run here?), \textbf{economic} (cost within budget? expected gain?), \textbf{operational} (can Mama Caro and Blessing use it daily?), \textbf{schedule} (can it be in place quickly?). Conclude with a feasibility verdict.
\item \textbf{Report and presentation.} Write a one-page needs analysis report with the headings: \emph{Context, Problem statement, Identified needs, Proposed IT solution, Feasibility, Expected benefits}. Prepare a 5-minute presentation in which each team member speaks.
\item \textbf{Class vote and reflection.} After all presentations, vote for the most realistic and impactful proposal (you may not vote for your own team). In two sentences, explain what made the winning proposal convincing.
\end{enumerate}

\courssub{Model Answer}

\begin{exoresolu}[Complete Needs Analysis for Mama Caro's Store]
\textbf{Statement:} Using the situation and observation data above, carry out the full needs analysis: list problems, design an interview guide, classify needs, match IT solutions, evaluate feasibility of the priority proposal, and draft the one-page report.

\tcblower

\textbf{1. Problems identified and their consequences.}

\begin{center}
\begin{tabularx}{\linewidth}{|l|X|X|}\hline
\rowcolor{popblueL} \textbf{No} & \textbf{Problem} & \textbf{Consequence} \\ \hline
1 & Paper records, torn pages, forgotten credit entries & Debts denied by customers; direct loss of money. \\ \hline
2 & No stock list, no low-stock alert & Stock-outs (three cartons of oil); two days of lost sales. \\ \hline
3 & MoMo/OM float mixed with shop cash & Impossible to know real profit; risk of errors and fraud. \\ \hline
4 & Cash kept in a box, no camera, no record backup & High loss in case of theft or fire. \\ \hline
5 & No advertising, no customer contacts & Customers attracted by competitors on WhatsApp. \\ \hline
6 & No monthly figures & Mama Caro cannot tell where the money goes; poor decisions. \\ \hline
\end{tabularx}
\end{center}

\textbf{2. Interview guide (six questions).}

\begin{enumerate}
\item How many times per month, on average, do you discover that an item has finished without warning? (frequency of stock-outs)
\item Approximately how much money is currently owed to you by credit customers? (size of the debt problem)
\item How do you separate the Mobile Money float from the shop's cash at the end of the day? (current payment practice)
\item If the exercise book were stolen or burnt, what information would you lose? (awareness of backup risk)
\item Would you and Blessing be willing to spend 15 minutes per day entering sales into a phone application? (operational willingness)
\item What is the first improvement you would like to see within three months? (owner's own priority)
\end{enumerate}

\textbf{3. Classification of needs.}

\begin{center}
\begin{tabularx}{\linewidth}{|l|X|X|}\hline
\rowcolor{popblueL} \textbf{Need category} & \textbf{Problem observed} & \textbf{Evidence from the field} \\ \hline
Record keeping and accounting & Sales and debts on paper, entries forgotten & Torn exercise book, loose papers, denied debts. \\ \hline
Stock management & No stock list, no alert & Oil stock-out discovered too late. \\ \hline
Payments and separation of funds & MoMo/OM float mixed with shop cash & Single book for all money movements. \\ \hline
Security & Cash box, no camera, no backup & Break-in at the neighbouring shop. \\ \hline
Marketing and customer relations & No advertising, no contact list & Competitors active on WhatsApp status. \\ \hline
\end{tabularx}
\end{center}

\textbf{4. Matching needs to IT solutions.}

\begin{center}
\begin{tabularx}{\linewidth}{|l|X|X|}\hline
\rowcolor{popblueL} \textbf{Need} & \textbf{Proposed IT solution} & \textbf{Justification} \\ \hline
Record keeping & A simple sales/debt notebook app or spreadsheet on Mama Caro's Android phone & Free or cheap, uses existing equipment, searchable and cannot be torn. \\ \hline
Stock management & Stock sheet (spreadsheet) with a minimum-level column checked weekly & No extra hardware; a simple formula flags items below the alert level. \\ \hline
Separation of funds & Two dedicated MoMo/OM agent lines plus a daily reconciliation sheet & Separates float from shop cash; daily totals can be checked against agent statements. \\ \hline
Security & Automatic cloud backup of the phone records; later a low-cost camera & Backup protects against theft/fire at almost zero cost; camera when budget allows. \\ \hline
Marketing & WhatsApp Business account with catalogue and status posts & Free, already familiar, reaches customers directly. \\ \hline
\end{tabularx}
\end{center}

\textbf{5. Priority proposal and feasibility checklist.}

Priority proposal: \textbf{a digital sales, debt and stock record system on Mama Caro's smartphone} (spreadsheet or shop-keeping app), with automatic cloud backup.

\begin{center}
\begin{tabularx}{\linewidth}{|l|X|X|}\hline
\rowcolor{popblueL} \textbf{Dimension} & \textbf{Question} & \textbf{Answer for this shop} \\ \hline
Technical & Can the technology run here? & Yes: Android phone available, mobile data coverage in Bamenda, generator covers power cuts. \\ \hline
Economic & Is it within budget? & Yes: app/spreadsheet free or under 10\,000 FCFA; total far below 150\,000 FCFA; recovered debts and avoided stock-outs quickly exceed the cost. \\ \hline
Operational & Can the users manage it daily? & Yes with conditions: Mama Caro already uses a smartphone; Blessing needs a short training and a fixed 15-minute daily routine. \\ \hline
Schedule & Can it be in place quickly? & Yes: installation, setup and training feasible within one to two weeks. \\ \hline
\end{tabularx}
\end{center}

\textbf{Verdict:} the proposal is \cle{feasible} on all four dimensions and is therefore retained as the priority.

\textbf{6. Extract of the needs analysis report.}

\begin{quote}
\textbf{Context.} Mama Caro's provision store (Mile 4, Bamenda), 9 years of activity, one employee, smartphone available, budget 150\,000 FCFA.\par
\textbf{Problem statement.} The business loses money and sales because records, stock and payments are managed on paper with no alerts, no separation of funds and no backup.\par
\textbf{Identified needs.} Reliable record keeping; stock monitoring; separation of MoMo/OM float; data security and backup; customer marketing.\par
\textbf{Proposed IT solution.} Digital sales, debt and stock records on the owner's smartphone with cloud backup, plus a WhatsApp Business catalogue.\par
\textbf{Feasibility.} Positive on the technical, economic, operational and schedule dimensions.\par
\textbf{Expected benefits.} Fewer denied debts, no unnoticed stock-outs, clear monthly profit, records safe from theft or fire, and new customers reached through WhatsApp.
\end{quote}

\textbf{7. Vote and reflection.} The winning proposal is usually the one that is cheap, uses equipment the owner already has, and solves the most painful problem first. A good reflection sentence would be: \emph{``The winning team convinced us because their solution costs almost nothing, starts from the owner's own smartphone, and attacks the problem that loses the most money --- the forgotten debts.''}

\end{exoresolu}

\begin{aretenir}[What this activity teaches]
A good IT needs analysis always follows the same path: \cle{observe} $\rightarrow$ \cle{interview} $\rightarrow$ \cle{classify} $\rightarrow$ \cle{match} $\rightarrow$ \cle{check feasibility}. The best solution is not the most technological one, but the one that is \textbf{useful, usable, affordable and sustainable} for that particular business.
\end{aretenir}

\begin{exercice}[Extension: Your Turn]
Choose a small business near your home or school (a hair salon, a call box, a food vendor, a tailor). Carry out the same 5-step analysis in 30 minutes. Write a half-page report and present it in 3 minutes. Compare your findings with Mama Caro's case: which needs are universal, which are specific?
\end{exercice}

\newpage

\coursec{Methods and worked examples}

\courssub{Methods and Worked Examples}

\begin{methode}[Analysing a Business Context to Extract Core Needs]
\textbf{Objective:} Move from a vague business description to a structured list of explicit and implicit needs.
\begin{enumerate}
    \item \textbf{Identify the stakeholders:} Owners, managers, employees, customers, suppliers, regulators (e.g., MINCOMMERCE, ANTIC).
    \item \textbf{Map the value chain:} Input $\to$ Processes $\to$ Output $\to$ Delivery $\to$ After-sales. Ask ``What happens at each stage?''.
    \item \textbf{Classify needs using PIECES:}
    \begin{itemize}
        \item \textbf{P}erformance (speed, throughput)
        \item \textbf{I}nformation (accuracy, timeliness, reporting)
        \item \textbf{E}conomy (cost reduction, resource optimisation)
        \item \textbf{C}ontrol (security, audit trail, compliance)
        \item \textbf{E}fficiency (automation, workflow)
        \item \textbf{S}ervice (customer experience, reliability)
    \end{itemize}
    \item \textbf{Distinguish symptoms from root causes:} ``Stockouts'' is a symptom; ``No real-time inventory tracking'' is a root cause.
    \item \textbf{Document each need as:} ``The business needs to \emph{<action>} so that \emph{<benefit>}.''
\end{enumerate}
\emph{Reflex:} Always validate needs with at least two stakeholders to avoid ``vocal minority'' bias.
\end{methode}

\begin{exoresolu}[Case Study: ``Mami Wata'' Restaurant in Douala]
\textbf{Statement:} ``Mami Wata'' is a popular mid-size restaurant in Akwa, Douala. The owner, Madame Ngono, complains: ``We lose money every day. Customers wait too long for bills, waiters forget orders, and we often run out of popular dishes while ingredients for unpopular ones rot in the fridge. My accountant spends three days a month reconciling cash with the notebook.''
\begin{enumerate}
    \item Identify four stakeholders and one specific pain point for each.
    \item Using the PIECES framework, classify the following needs: (a) Reduce billing time from 10 min to 2 min. (b) Know exactly which dishes are profitable. (c) Prevent waiters from giving free meals to friends. (d) Allow customers to order via QR code at the table.
    \item Rewrite the owner's complaint into three structured need statements using the format ``The business needs to \emph{<action>} so that \emph{<benefit>}.''
\end{enumerate}
\tcblower
\textbf{Detailed Solution:}
\begin{enumerate}
    \item \textbf{Stakeholders \& Pain Points:}
    \begin{center}
    \begin{tabularx}{\linewidth}{|l|X|}\hline
    \rowcolor{popblueL} \textbf{Stakeholder} & \textbf{Pain Point} \\ \hline
    Owner (Madame Ngono) & Low profitability, time wasted on manual reconciliation. \\ \hline
    Waiters & Stress from forgotten orders, conflicts with kitchen, no tips due to slow service. \\ \hline
    Kitchen Staff & Chaos from verbal orders, ingredient waste, unable to prepare efficiently. \\ \hline
    Customers & Long waits for food and bill, errors in orders, no digital payment option (Mobile Money). \\ \hline
    Accountant (External) & Three days of manual data entry from error-prone notebooks. \\ \hline
    \end{tabularx}
    \end{center}
    \item \textbf{PIECES Classification:}
    \begin{itemize}
        \item (a) Reduce billing time $\to$ \textbf{E}fficiency (workflow automation) and \textbf{S}ervice (customer wait time).
        \item (b) Know profitable dishes $\to$ \textbf{I}nformation (accurate, timely reporting) and \textbf{E}conomy (cost control).
        \item (c) Prevent free meals $\to$ \textbf{C}ontrol (security, audit trail, fraud prevention).
        \item (d) QR code ordering $\to$ \textbf{S}ervice (customer experience) and \textbf{E}fficiency (reduces waiter trips).
    \end{itemize}
    \item \textbf{Structured Need Statements:}
    \begin{enumerate}
        \item The business needs to \emph{automate order taking and billing via a POS system} so that \emph{service time is reduced, errors are eliminated, and each transaction is recorded instantly}.
        \item The business needs to \emph{track ingredient consumption per dish sold in real time} so that \emph{stockouts of popular items are avoided and waste from perishable ingredients is minimised}.
        \item The business needs to \emph{generate daily sales and reconciliation reports automatically} so that \emph{the accountant spends hours instead of days on month-end closing and the owner has real-time visibility on profitability}.
    \end{enumerate}
\end{enumerate}
\end{exoresolu}

\begin{methode}[Conducting Effective Requirements Elicitation]
\textbf{Objective:} Gather complete, accurate, and prioritised requirements using complementary techniques.
\begin{enumerate}
    \item \textbf{Prepare a Stakeholder Map:} Power/Interest grid (High Power/High Interest = Manage Closely; Low Power/Low Interest = Monitor).
    \item \textbf{Select Techniques (Triangulation):}
    \begin{itemize}
        \item \textbf{Interviews (Structured/Semi-structured):} Deep dive, build rapport. Use open questions: ``Walk me through a typical day...''
        \item \textbf{Observation / Job Shadowing:} Reveals tacit knowledge, workarounds, actual vs. documented process. \emph{Crucial for ICT:} Watch how they use current tools (Excel, WhatsApp, paper).
        \item \textbf{Document Analysis:} Invoices, receipts, ledgers, regulatory forms (business registration, tax returns), existing reports.
        \item \textbf{Workshops / JAD (Joint Application Design):} Resolve conflicts, prioritise, model processes together.
        \item \textbf{Questionnaires / Surveys:} Quantify frequency and satisfaction (large groups such as customers).
        \item \textbf{Prototyping / Mock-ups:} ``Show, don't just tell.'' Low-fidelity wireframes (paper) $\to$ High-fidelity (Figma).
    \end{itemize}
    \item \textbf{Record Requirements as User Stories:} ``As a \emph{<role>}, I want to \emph{<action>}, so that \emph{<benefit>}.'' Add \textbf{Acceptance Criteria} (Given/When/Then).
    \item \textbf{Prioritise using MoSCoW:} \textbf{M}ust have, \textbf{S}hould have, \textbf{C}ould have, \textbf{W}on't have (this phase).
    \item \textbf{Validate \& Sign-off:} Walkthrough with stakeholders; traceability matrix (Requirement $\leftrightarrow$ Source $\leftrightarrow$ Test Case).
\end{enumerate}
\emph{Warning:} Never rely on a single technique. Interviews reveal \emph{what people say}; observation reveals \emph{what people do}.
\end{methode}

\begin{exoresolu}[Elicitation Plan for a Mobile Money Agent Network]
\textbf{Statement:} You are hired by a fintech startup ``PayFast'' to design a management dashboard for their network of 150 Mobile Money agents across the Littoral and Centre regions. Agents currently use paper registers and USSD menus. The startup wants to reduce fraud, monitor liquidity in real time, and automate commission calculations.
\begin{enumerate}
    \item Draw a Power/Interest grid placing: (a) PayFast CEO, (b) Regional Field Supervisor, (c) Individual Agent (rural), (d) Individual Agent (urban, high volume), (e) BEAC Regulator, (f) MTN/Orange API Team.
    \item Propose a triangulation plan (3 techniques) to elicit requirements from the agents, justifying why each is needed given the context (literacy levels, connectivity, trust).
    \item Write two User Stories with Acceptance Criteria for: (i) Real-time liquidity monitoring, (ii) Automated commission calculation.
\end{enumerate}
\tcblower
\textbf{Detailed Solution:}
\begin{enumerate}
    \item \textbf{Power/Interest Grid:}
    \begin{center}
    \begin{tabularx}{\linewidth}{|l|l|X|}\hline
    \rowcolor{popblueL} \textbf{Stakeholder} & \textbf{Power} & \textbf{Interest (Strategy)} \\ \hline
    PayFast CEO & High & High (Manage Closely) \\ \hline
    Regional Field Supervisor & High & High (Manage Closely) \\ \hline
    Rural Agent & Low & High (Keep Informed) \\ \hline
    Urban High-Volume Agent & Medium & High (Keep Satisfied) \\ \hline
    BEAC Regulator & High & Medium (Keep Satisfied -- Compliance) \\ \hline
    MTN/Orange API Team & High & Low (Monitor -- Technical dependency) \\ \hline
    \end{tabularx}
    \end{center}
    \item \textbf{Triangulation Plan for Agents:}
    \begin{enumerate}
        \item \textbf{Observation / Job Shadowing (2 days per region):} \emph{Why:} Many agents have low formal literacy; they use USSD codes by heart. Watching them reveals error-prone steps (e.g., mistyping a long USSD string), how they handle float (cash vs e-money), and trust issues with supervisors. Interviews alone would miss these tacit workflows.
        \item \textbf{Semi-structured Interviews (with Supervisors translating if needed):} \emph{Why:} To capture pain points in their own words (e.g., ``I don't know my commission until the 5th of next month''), understand local constraints (power cuts, network downtime), and build trust for adoption.
        \item \textbf{Document Analysis (Paper Registers + USSD Logs + Bank Statements):} \emph{Why:} Quantify transaction volumes, error rates, reconciliation discrepancies. Provides baseline data for dashboard KPIs (e.g., average float holding period).
    \end{enumerate}
    \item \textbf{User Stories:}
    \begin{enumerate}
        \item \textbf{Story:} As a \emph{Regional Field Supervisor}, I want to \emph{see the real-time e-float and cash float balance of every agent on a map}, so that \emph{I can redeploy liquidity before an agent runs dry or holds excess cash}.
        \begin{itemize}
            \item \textbf{Acceptance Criteria:}
            \begin{itemize}
                \item Given the dashboard is open, When an agent completes a deposit/withdrawal via USSD/API, Then their balance updates within 30 seconds.
                \item Given an agent's e-float $< 10\%$ of daily average volume, Then the agent pin turns \textbf{Red} and an SMS alert is sent to the supervisor.
                \item Given a network outage, Then the last known balance is shown with a ``Stale Data'' timestamp.
            \end{itemize}
        \end{itemize}
        \item \textbf{Story:} As an \emph{Agent}, I want to \emph{see my daily, weekly, and monthly commission breakdown per transaction type (deposit, withdrawal, bill pay)}, so that \emph{I can verify my earnings instantly and plan my float}.
        \begin{itemize}
            \item \textbf{Acceptance Criteria:}
            \begin{itemize}
                \item Given the agent logs into the mobile app/USSD menu, When they select ``My Commissions'', Then an itemised list appears: Date, Type, Amount, Rate, Commission.
                \item Given month-end, When the supervisor approves, Then a PDF statement is generated and sent via WhatsApp/Email.
                \item Commission formula (example): Deposit $0.5\%$, Withdrawal $1.2\%$, Bill Pay $50$ FCFA flat. Rates configurable by admin without code change.
            \end{itemize}
        \end{itemize}
    \end{enumerate}
\end{enumerate}
\end{exoresolu}

\begin{methode}[Modelling Business Processes (AS-IS vs TO-BE)]
\textbf{Objective:} Visualise current workflows to identify bottlenecks, redundancies, and automation opportunities.
\begin{enumerate}
    \item \textbf{Choose Notation:} BPMN 2.0 (Business Process Model and Notation) is standard. Key elements:
    \begin{itemize}
        \item \textbf{Pool/Lane:} Participant (e.g., Customer, Waiter, Kitchen, Cashier, System).
        \item \textbf{Start Event (thin circle):} Trigger (Message, Timer, Signal).
        \item \textbf{Task (rounded rectangle):} Work unit (User Task, Service Task/Script, Send/Receive Task).
        \item \textbf{Gateway (diamond):} Decision (Exclusive XOR, Parallel AND, Inclusive OR).
        \item \textbf{Flow:} Sequence (solid arrow), Message (dashed arrow).
        \item \textbf{End Event (thick circle):} Outcome.
    \end{itemize}
    \item \textbf{Model AS-IS (Current State):} Interview + Observation. Capture \emph{exceptions} (``What if stock out?''). Measure: Cycle time, Wait time, Handoffs, Error rate.
    \item \textbf{Analyse Waste (Lean / TIMWOODS):} \textbf{T}ransport, \textbf{I}nventory, \textbf{M}otion, \textbf{W}aiting, \textbf{O}ver-processing, \textbf{O}ver-production, \textbf{D}efects, \textbf{S}kills underutilised.
    \item \textbf{Design TO-BE (Future State):} Apply IT to eliminate waste: automate Service Tasks, replace manual Message Flows with API calls, add Decision Gateways with business rules, introduce parallel paths.
    \item \textbf{Validate TO-BE:} Walkthrough with stakeholders; simulate with sample data; check regulatory compliance (e.g., OHADA accounting, BEAC KYC).
\end{enumerate}
\emph{Exam Tip:} In the GCE, you may be asked to draw a process fragment. Practise drawing: Start $\to$ User Task $\to$ XOR Gateway $\to$ [Yes] Service Task $\to$ [No] User Task $\to$ End.
\end{methode}

\begin{exoresolu}[Process Modelling: School Fees Collection in a Cameroonian Lycée]
\textbf{Statement:} At a lycée in Yaoundé, fees collection is manual. Process:
\begin{enumerate}
    \item Parent arrives at Bursar's office, queues.
    \item Bursar checks paper ledger for student balance.
    \item Parent pays cash. Bursar writes receipt (carbon copy), updates ledger manually.
    \item Parent takes original receipt to Class Teacher for signature (proof of payment).
    \item Class Teacher signs, parent keeps copy.
    \item End of day: Bursar counts cash, fills deposit slip, walks to bank. Next day: Bank confirms deposit.
\end{enumerate}
Problems: Queues (2h), lost receipts, ledger errors, theft risk, delayed bank reconciliation.
\begin{enumerate}
    \item Describe the AS-IS BPMN diagram (lanes/tasks/gateways in text form).
    \item Identify 3 TIMWOODS wastes with evidence from the scenario.
    \item Design the TO-BE process using a Mobile Money (MoMo) payment integration and a simple School Management System (SMS). Describe the new flow and the TO-BE BPMN fragment for the ``Parent pays via MoMo'' path.
    \item List 2 non-functional requirements critical for this TO-BE system in the Cameroonian context.
\end{enumerate}
\tcblower
\textbf{Detailed Solution:}
\begin{enumerate}
    \item \textbf{AS-IS BPMN (Textual Description):}
    \begin{itemize}
        \item \textbf{Pools/Lanes:} Parent, Bursar, Class Teacher, Bank.
        \item \textbf{Start:} Parent arrives (Message Event).
        \item \textbf{Tasks:}
        \begin{itemize}
            \item Parent: Queue (Wait) $\to$ Pay Cash (User Task) $\to$ Collect Receipt $\to$ Walk to Teacher $\to$ Get Signature.
            \item Bursar: Search Ledger $\to$ Write Receipt $\to$ Update Ledger $\to$ End-of-Day: Count Cash $\to$ Fill Deposit Slip $\to$ Walk to Bank.
            \item Class Teacher: Sign Receipt.
            \item Bank: Process Deposit (external Service Task).
        \end{itemize}
        \item \textbf{Gateways:} XOR at Bursar: ``Student found in ledger?'' [No $\to$ Manual search/Principal office]. XOR at Teacher: ``Receipt valid?'' [No $\to$ Reject].
        \item \textbf{End:} Parent has signed receipt; Bank deposit confirmed (next day).
    \end{itemize}
    \item \textbf{TIMWOODS Wastes:}
    \begin{enumerate}
        \item \textbf{Waiting:} Parents queue 2h; Bursar waits for bank confirmation the next day.
        \item \textbf{Motion:} Parent walks Bursar $\to$ Teacher; Bursar walks physically to the Bank.
        \item \textbf{Defects:} Lost receipts, ledger errors (wrong amount, wrong student), cash theft risk.
        \item \emph{(Bonus) Over-processing:} Carbon copy receipt + ledger entry + teacher signature = triple recording.
    \end{enumerate}
    \item \textbf{TO-BE Process (MoMo + SMS):}
    \textbf{New Flow:}
    \begin{enumerate}
        \item Parent pays via a MoMo merchant/paybill code, entering the Student ID and amount (USSD or MoMo App).
        \item The School Management System validates the Student ID, updates the balance, and sends instant SMS confirmation to Parent + Bursar + Class Teacher.
        \item Bursar dashboard updates in real time (no manual ledger).
        \item Auto-reconciliation: MoMo transaction report $\to$ SMS matches $\to$ Bank settlement file generated automatically.
        \item Class Teacher sees ``Paid'' status instantly on their phone/portal -- no physical signature needed.
    \end{enumerate}
    \textbf{TO-BE BPMN Fragment (Parent Pays via MoMo):}
    \begin{itemize}
        \item \textbf{Lanes:} Parent, MoMo System (External), School Management System, Bursar (Monitor), Class Teacher (Monitor).
        \item \textbf{Start Event (Message):} Parent initiates USSD/App payment.
        \item \textbf{Service Task (MoMo System):} Validate PIN, Check Float, Debit Parent, Credit School account.
        \item \textbf{Message Flow:} MoMo $\to$ SMS: Transaction Callback (StudentID, Amount, Ref, Timestamp).
        \item \textbf{Service Task (SMS):} Verify StudentID, Update Balance, Generate Digital Receipt, Trigger Notifications.
        \item \textbf{Parallel Gateway (AND):} Split to three Send Tasks: Send SMS to Parent (Receipt Ref); Update Bursar Dashboard (real time); Update Teacher Portal (Paid flag).
        \item \textbf{End Event:} Payment Confirmed.
    \end{itemize}
    \item \textbf{Critical Non-Functional Requirements (Cameroon Context):}
    \begin{enumerate}
        \item \textbf{Offline Resilience / USSD Fallback:} Internet is unreliable in many areas. The system \emph{must} support USSD initiation (no smartphone/data needed) and queue callbacks for retry if the school's connection is down, syncing when connectivity returns.
        \item \textbf{Audit Trail \& OHADA Compliance:} Every transaction must generate an immutable log (timestamp, user, action, before/after values) exportable for OHADA-compliant financial statements and audits. No ``delete'' function for financial records -- only ``reverse'' with authorisation.
    \end{enumerate}
\end{enumerate}
\end{exoresolu}

\begin{methode}[Mapping Business Needs to IT Solution Categories]
\textbf{Objective:} Classify needs and match them to appropriate IT solution types (Buy, Build, Integrate, Cloud).
\begin{enumerate}
    \item \textbf{Categorise each need:}
    \begin{itemize}
        \item \textbf{Operational / Transactional:} High volume, structured, real-time (POS, ERP core, Mobile Money API). $\to$ \textbf{Buy/Configure} (Odoo, Sage, custom POS) or \textbf{Build} if unique IP.
        \item \textbf{Analytical / Decision Support:} Aggregated, historical, ad-hoc queries (Dashboards, BI). $\to$ \textbf{Buy} (Power BI, Metabase, Superset) + Data Warehouse.
        \item \textbf{Collaborative / Communication:} Unstructured, human-centric (Email, WhatsApp Business, Ticketing). $\to$ \textbf{Buy/SaaS} (Zoho, Freshdesk, Microsoft 365).
        \item \textbf{Customer-Facing / Digital Channel:} Web/App, Public API, E-commerce. $\to$ \textbf{Build} (React/Flutter + Node/Django) or \textbf{Low-Code} for speed.
        \item \textbf{Infrastructure / Security:} Hosting, Backup, Firewall, Identity (SSO/MFA). $\to$ \textbf{Cloud/IaaS/PaaS} (AWS, Azure, OVH, local MTN Cloud) + Managed Services.
    \end{itemize}
    \item \textbf{Apply a Decision Matrix (Weighted Scoring):} Score each option 1--5 per criterion, multiply by the weight, and sum.
    \begin{center}
    \begin{tabularx}{\linewidth}{|X|c|c|c|c|}\hline
    \rowcolor{popblueL} \textbf{Criterion (Weight)} & \textbf{Buy SaaS} & \textbf{Custom Build} & \textbf{Low-Code} & \textbf{Open Source} \\ \hline
    Time to Market (30\%) & 5 & 1 & 4 & 2 \\ \hline
    Total Cost 3yr (25\%) & 3 & 2 & 4 & 5 \\ \hline
    Customisation Fit (20\%) & 2 & 5 & 4 & 4 \\ \hline
    Data Sovereignty (15\%) & 1 & 5 & 2 & 5 \\ \hline
    Local Support (10\%) & 3 & 4 & 3 & 2 \\ \hline
    \textbf{Weighted Score} & \textbf{2.95} & \textbf{2.90} & \textbf{3.60} & \textbf{3.55} \\ \hline
    \end{tabularx}
    \end{center}
    \emph{Example calculation (Buy SaaS):} $5(0.30)+3(0.25)+2(0.20)+1(0.15)+3(0.10)=1.5+0.75+0.4+0.15+0.3=3.10$. Recompute each column the same way; the highest score wins, subject to constraints below.
    \item \textbf{Consider Cameroon Constraints:}
    \begin{itemize}
        \item \textbf{Connectivity:} Favour offline-first, USSD, Progressive Web Apps (PWA).
        \item \textbf{Power:} Lightweight clients, battery-efficient mobile apps.
        \item \textbf{Skills:} Low-code enables local developers; SaaS reduces operations burden.
        \item \textbf{Regulation:} ANTIC guidance on data protection and localisation $\to$ prefer local hosting or hybrid.
        \item \textbf{Payment:} Must integrate Mobile Money (MTN MoMo API, Orange Money API) natively.
    \end{itemize}
    \item \textbf{Output:} Solution Architecture Diagram (Context + Containers + Components).
\end{enumerate}
\end{methode}

\begin{exoresolu}[Solution Mapping for ``AgroConnect'' -- Linking Farmers to Buyers]
\textbf{Statement:} ``AgroConnect'' wants to build a platform connecting smallholder cocoa/coffee farmers in the South-West Region to exporters in Douala. Key needs:
\begin{enumerate}
    \item Farmers (low literacy, basic phones, poor internet) need to know daily market prices and sell produce.
    \item Field agents (motorcycles, smartphones) visit farms, grade quality, record weight, issue digital receipts.
    \item Exporters need real-time aggregate supply view, quality reports, and traceability.
    \item Payments: Farmers paid via Mobile Money; Exporters pay via Bank Transfer.
    \item Platform must work offline for agents; sync when online.
\end{enumerate}
\begin{enumerate}
    \item Classify each need (1--5) into the five solution categories (Operational, Analytical, Collaborative, Customer-Facing, Infrastructure).
    \item Propose a technology stack (Frontend, Backend, Database, Offline Sync, Payments, Hosting) justifying each choice for the Cameroonian context.
    \item Draw a high-level Context Diagram showing: Farmers, Agents, Exporters, Admins, MoMo API, Bank API, SMS Gateway, Platform.
\end{enumerate}
\tcblower
\textbf{Detailed Solution:}
\begin{enumerate}
    \item \textbf{Classification:}
    \begin{center}
    \begin{tabularx}{\linewidth}{|c|X|l|}\hline
    \rowcolor{popblueL} \textbf{Need} & \textbf{Description} & \textbf{Category} \\ \hline
    1 & Farmer price info \& sell intent (USSD/Voice) & Customer-Facing (Digital Channel) \\ \hline
    2 & Agent: Grade, weigh, digital receipt (Offline Mobile App) & Operational (Transactional) \\ \hline
    3 & Exporter: Supply dashboard, quality, traceability & Analytical (Decision Support) \\ \hline
    4 & Payments: MoMo to farmers, Bank to exporters & Operational + Infrastructure (Integration) \\ \hline
    5 & Offline-first sync for agents & Infrastructure (Platform Capability) \\ \hline
    \end{tabularx}
    \end{center}
    \item \textbf{Technology Stack Justification:}
    \begin{center}
    \begin{tabularx}{\linewidth}{|l|l|X|}\hline
    \rowcolor{popblueL} \textbf{Layer} & \textbf{Choice} & \textbf{Justification (Cameroon Context)} \\ \hline
    Farmer Interface & USSD + IVR (Voice) & Works on any phone, no internet, no literacy barrier (local languages). \\ \hline
    Agent App & Flutter (Android) + SQLite local DB & Single codebase; local DB for offline work; sync engine pushes data when the agent reaches 3G/4G coverage. \\ \hline
    Exporter Portal & React (PWA) & Responsive, works on modest laptops; PWA allows ``install'' on desktop. \\ \hline
    Backend API & Node.js or Python (FastAPI) & Local developer talent available; FastAPI auto-generates API documentation for exporter integration. \\ \hline
    Database & PostgreSQL + PostGIS & ACID for transactions; PostGIS for farm geo-location (traceability mapping). \\ \hline
    Offline Sync & Sync service (SQLite $\leftrightarrow$ Postgres) & Handles conflict resolution and schema migrations for offline-first mobile. \\ \hline
    Payments & MTN MoMo API + Orange Money API + Bank transfer & Covers the large majority of the mobile money market; bank API for high-value exporter settlements. \\ \hline
    SMS/Notifications & SMS aggregator (e.g., Infobip, Termii) & Delivery receipts critical for digital receipts; supports Sender ID ``AgroConnect''. \\ \hline
    Hosting & Local/regional cloud (e.g., MTN Cloud) + EU backup & Data sovereignty concerns (ANTIC); low latency for Douala exporters. \\ \hline
    Auth & Keycloak (self-hosted) & SSO, MFA, Role-Based Access Control (Farmer, Agent, Exporter, Admin); no vendor lock-in. \\ \hline
    \end{tabularx}
    \end{center}
    \item \textbf{Context Diagram (Level 1):}
    \begin{popfigure}
    \begin{center}
    \begin{tikzpicture}[every node/.style={font=\small}]
        \node[draw, rounded corners, fill=popblueL, minimum width=3.5cm, minimum height=1.6cm, align=center] (Platform) at (0,0) {\textbf{AgroConnect}\\ \textbf{Platform}};
        \node[draw, circle, fill=popgreenL, align=center, minimum size=1.6cm] (Farmer) at (-5,0) {\textbf{Farmer}\\ USSD/IVR};
        \node[draw, circle, fill=poporange, align=center, minimum size=1.6cm] (Agent) at (-3.5,3) {\textbf{Agent}\\ App};
        \node[draw, circle, fill=poppurple, align=center, minimum size=1.6cm] (Exporter) at (3.5,3) {\textbf{Exporter}\\ PWA};
        \node[draw, circle, fill=poppink, align=center, minimum size=1.6cm] (Admin) at (5,0) {\textbf{Admin}\\ PWA};
        \node[draw, rectangle, dashed, fill=popblueL, align=center] (MoMo) at (-3.5,-3) {\textbf{MoMo API}};
        \node[draw, rectangle, dashed, fill=popblueL, align=center] (Bank) at (3.5,-3) {\textbf{Bank API}};
        \node[draw, rectangle, dashed, fill=popblueL, align=center] (SMS) at (0,-3.2) {\textbf{SMS Gateway}};
        \draw[->, thick] (Farmer) -- node[above, sloped] {USSD/Voice} (Platform);
        \draw[<->, thick] (Agent) -- node[above, sloped] {HTTPS/Sync} (Platform);
        \draw[<->, thick] (Exporter) -- node[above, sloped] {HTTPS/REST} (Platform);
        \draw[<->, thick] (Admin) -- node[above] {HTTPS} (Platform);
        \draw[<->, thick, dashed] (Platform) -- node[below, sloped] {Disburse/Collect} (MoMo);
        \draw[<->, thick, dashed] (Platform) -- node[below, sloped] {Settlement} (Bank);
        \draw[->, thick, dashed] (Platform) -- node[right] {Receipts} (SMS);
    \end{tikzpicture}
    \end{center}
    \legende{Context diagram of the AgroConnect platform. Solid circles: human actors; dashed rectangles: external systems.}
    \end{popfigure}
\end{enumerate}
\end{exoresolu}

\begin{methode}[Feasibility Analysis: TELOS]
\textbf{Objective:} Determine if the proposed IT solution is viable before committing resources. TELOS = \textbf{T}echnical, \textbf{E}conomic, \textbf{L}egal, \textbf{O}perational, \textbf{S}chedule.
\begin{enumerate}
    \item \textbf{Technical Feasibility:} Can we build it with available technology and skills?
        \begin{itemize}
            \item Proven tech stack? (Avoid ``bleeding edge''.)
            \item Integration complexity? (Are APIs available? MoMo API sandbox access?)
            \item Scalability/performance targets achievable? (e.g., peak of many concurrent USSD sessions.)
            \item Offline sync conflict-resolution strategy defined?
            \item Team skills match? (Mobile, backend, database, DevOps.)
        \end{itemize}
    \item \textbf{Economic Feasibility (Cost-Benefit Analysis):}
        \begin{itemize}
            \item \textbf{CAPEX (one-off):} Development, infrastructure setup, licences.
            \item \textbf{OPEX (recurring):} Cloud hosting, SMS, API fees, support, maintenance (often estimated at 15--25\% of CAPEX per year).
            \item \textbf{Benefits:} New revenue (commissions, subscriptions), cost savings, productivity gains.
            \item \textbf{Indicators:} Net Present Value (NPV), Return on Investment (ROI), Payback Period. A short payback (e.g., under 18--24 months) is generally a good sign for a startup.
        \end{itemize}
    \item \textbf{Legal/Regulatory Feasibility (Cameroon):}
        \begin{itemize}
            \item \textbf{ANTIC:} Personal data protection obligations, cybersecurity standards for electronic services.
            \item \textbf{BEAC/COBAC:} CEMAC regulations on payment services and e-money; KYC/AML (Know Your Customer / Anti-Money Laundering) duties.
            \item \textbf{MINCOMMERCE / Trade Register:} Business registration; e-commerce declarations where applicable.
            \item \textbf{OHADA:} Accounting compliance (SYSCOHADA chart of accounts).
            \item \textbf{Intellectual Property:} OAPI registration for brand and software protection.
        \end{itemize}
    \item \textbf{Operational Feasibility:} Will users adopt it? Is there a change-management plan (training, incentives, helpdesk, local-language interfaces)?
    \item \textbf{Schedule Feasibility:} Can a Minimum Viable Product be delivered within the available time and budget? Break the work into monthly milestones ending with a pilot.
    \item \textbf{Go/No-Go Decision:} Document risks (technical, regulatory, market) and mitigations.
\end{enumerate}
\end{methode}

\begin{exoresolu}[Feasibility Study for ``EduPay'' -- School Fees Financing]
\textbf{Statement:} ``EduPay'' proposes a fintech platform allowing parents to pay school fees in monthly instalments via Mobile Money, while schools receive the full amount upfront (EduPay advances the funds). EduPay partners with 20 private schools in Yaoundé/Douala.
Key parameters: Average annual fee = 300,000 FCFA. Target: 5,000 students in Year 1. EduPay charges 1.5\%/month interest. Default rate estimated at 5\%. Cost of funds (borrowing from a bank) = 12\%/year. Operational costs (staff, cloud, SMS, marketing) = 50M FCFA/year. Development cost (MVP) = 40M FCFA.
\begin{enumerate}
    \item Perform a simplified Economic Feasibility (Year 1 projection). Calculate the Net Result.
    \item Identify 3 Legal/Regulatory hurdles specific to Cameroon and propose mitigations.
    \item Assess Operational Feasibility: Why might schools refuse? Why might parents default? Propose 2 mitigation strategies for each.
    \item Technical Risk: The Mobile Money Disbursement API (paying schools upfront) has a 2\% failure rate (timeouts). Design a reconciliation and retry mechanism (idempotency).
\end{enumerate}
\tcblower
\textbf{Detailed Solution:}
\begin{enumerate}
    \item \textbf{Economic Feasibility (Year 1):}
    \begin{align*}
    \text{Total Fee Volume} &= 5,000 \times 300,000 = 1,500,000,000 \text{ FCFA} \\
    \text{Average Loan Duration} &= 6 \text{ months (mid-year average)} \\
    \text{Interest Revenue} &= 1.5\text{B} \times 1.5\% \times 6 = 135,000,000 \text{ FCFA} \\
    \text{Default Loss (Principal)} &= 1.5\text{B} \times 5\% = 75,000,000 \text{ FCFA} \\
    \text{Cost of Funds} &= 1.5\text{B} \times 12\% \times 0.5 = 90,000,000 \text{ FCFA} \\
    \text{Gross Margin} &= 135\text{M} - 75\text{M} - 90\text{M} = -30,000,000 \text{ FCFA} \\
    \text{Operational Costs} &= 50,000,000 \text{ FCFA} \\
    \text{Amortisation (40M over 3 years)} &\approx 13,300,000 \text{ FCFA} \\
    \text{Net Result (Year 1)} &= -30\text{M} - 50\text{M} - 13.3\text{M} \approx \mathbf{-93.3 \text{ Million FCFA}}
    \end{align*}
    \textbf{Analysis:} A Year 1 loss is common (customer acquisition, default provisioning), but here the \emph{gross margin itself is negative}, which is structural. Break-even requires adjustments: lower default rate ($<3\%$), higher margin, cheaper funds, or greater scale. \emph{Not viable at the current parameters without changes.}
    \item \textbf{Legal/Regulatory Hurdles:}
    \begin{enumerate}
        \item \textbf{BEAC/COBAC regulation of credit activities:} EduPay is effectively lending money. This generally requires a licensed microfinance or payment institution status, or a partnership with one. \emph{Mitigation:} Partner with an existing licensed microfinance institution or bank, or seek the appropriate payment-institution authorisation.
        \item \textbf{ANTIC / personal data protection:} Processing minors' data (students) and parents' financial data requires a lawful basis, consent management, security measures, and respect of data-subject rights. \emph{Mitigation:} Privacy by Design; minimal data collection; local or compliant hosting; appoint a data-protection officer.
        \item \textbf{MINSEC / school partnership rules:} Private schools operate under MINSEC authorisation; financial partnerships must be contractually sound, and public schools cannot freely participate. \emph{Mitigation:} Legal vetting of school contracts; restrict the pilot to authorised private schools; use a pre-approved MoU template.
    \end{enumerate}
    \item \textbf{Operational Feasibility:}
    \begin{center}
    \begin{tabularx}{\linewidth}{|l|X|X|}\hline
    \rowcolor{popblueL} \textbf{Risk} & \textbf{Why} & \textbf{Mitigation} \\ \hline
    School Refusal & Fear of losing the direct relationship with parents; distrust of fintech; cash flow already managed. & 1. \textbf{Revenue share:} offer the school a small administration fee. 2. \textbf{Guarantee:} EduPay bears 100\% of default risk; the school receives the full fee on Day 1 regardless. \\ \hline
    Parent Default & Income volatility (informal sector); multiple simultaneous loans; no credit history. & 1. \textbf{Alternative scoring:} use Mobile Money transaction history and past school-payment behaviour. 2. \textbf{Flexible repayment:} align due dates with salary/harvest cycles; offer a supervised ``skip a month'' option. \\ \hline
    \end{tabularx}
    \end{center}
    \item \textbf{Technical: Idempotent Disbursement \& Reconciliation:}
    \begin{enumerate}
        \item \textbf{Idempotency Key:} Every disbursement request carries a unique key (e.g., built from school ID + term + batch ID). Re-sending the same request with the same key never pays twice.
        \item \textbf{State Machine:} \texttt{INITIATED} $\to$ \texttt{API\_SENT} $\to$ (\texttt{SUCCESS} | \texttt{TIMEOUT} | \texttt{FAILED}).
        \item \textbf{Retry Logic (exponential backoff):} On timeout, query the MoMo transaction-status API with the same key after 1 min, then 2, 5, 15 min. If \texttt{SUCCESS}, mark \texttt{COMPLETED}; if \texttt{NOT\_FOUND}, retry the disbursement (same key) up to 3 times.
        \item \textbf{Daily Reconciliation Job (e.g., 02:00):} Pull the MoMo settlement report; match by idempotency key against the internal ledger; flag exceptions (local SUCCESS vs MoMo FAILED; missed callbacks; duplicates) and produce a report for the finance team and auditors.
    \end{enumerate}
\end{enumerate}
\end{exoresolu}

\begin{methode}[Spotting and Validating IT Business Opportunities]
\textbf{Objective:} Move from ``I have an idea'' to ``This is a fundable, scalable venture'' using structured validation.
\begin{enumerate}
    \item \textbf{Problem-Solution Fit (Lean Canvas):} Fill a 1-page Lean Canvas:
    \begin{itemize}
        \item \textbf{Problem:} Top 3 pains (evidence: interviews, data).
        \item \textbf{Customer Segments:} Early adopters (who feels the pain most acutely?).
        \item \textbf{Unique Value Proposition (UVP):} ``We help \emph{<segment>} \emph{<achieve benefit>} by \emph{<differentiator>}.''
        \item \textbf{Solution:} Minimum feature set (MVP).
        \item \textbf{Channels:} How to reach customers (WhatsApp groups, radio, field agents, USSD).
        \item \textbf{Revenue Streams:} Transaction fee, subscription, freemium, data/ads.
        \item \textbf{Cost Structure:} CAC (Customer Acquisition Cost), COGS (API fees, SMS), operations.
        \item \textbf{Key Metrics:} Active users, retention, LTV/CAC $> 3$, payback $< 6$ months.
        \item \textbf{Unfair Advantage:} Exclusive partnership, proprietary data, licence, network effects.
    \end{itemize}
    \item \textbf{Market Sizing (TAM/SAM/SOM):}
    \begin{itemize}
        \item \textbf{TAM (Total Addressable Market):} All potential users $\times$ maximum price.
        \item \textbf{SAM (Serviceable Available Market):} The segment you can actually reach with your channels.
        \item \textbf{SOM (Serviceable Obtainable Market):} Realistic capture in Years 1--3.
    \end{itemize}
    \item \textbf{Validation Experiments (Build-Measure-Learn):}
    \begin{itemize}
        \item \textbf{Smoke Test:} Landing page + ``Join Waitlist'' button $\to$ measure sign-ups.
        \item \textbf{Concierge MVP:} Manually deliver the service (Excel + WhatsApp) for a few pilot clients $\to$ learn workflow, pricing, objections.
        \item \textbf{Pre-sell:} Collect Letters of Intent or small deposits before building.
        \item \textbf{Usability Test:} Paper prototype with about 5 users $\to$ task completion rate, time, errors.
    \end{itemize}
    \item \textbf{Pitch Deck Structure (10 slides):} Title $\to$ Problem $\to$ Solution $\to$ Market $\to$ Business Model $\to$ Traction $\to$ Team $\to$ Financials $\to$ Ask $\to$ Vision.
\end{enumerate}
\end{methode}

\begin{exoresolu}[Validating ``KmerJobs'' -- Blue-Collar Labour Marketplace]
\textbf{Statement:} You want to launch ``KmerJobs'', a platform connecting informal skilled workers (plumbers, electricians, carpenters, cleaners) in Douala with households/SMEs needing one-off jobs. Workers have smartphones (WhatsApp) but low formal education. Clients want trust, transparent pricing, and reliability.
\begin{enumerate}
    \item Complete a Lean Canvas for KmerJobs.
    \item Estimate TAM/SAM/SOM for Year 1 in Douala using these \emph{working assumptions} (illustrative, not official statistics): households $\approx$ 800k; SMEs $\approx$ 50k; informal skilled workers $\approx$ 200k; average job = 15,000 FCFA; platform fee = 10\%.
    \item Design a ``Concierge MVP'' experiment to validate the core risk: \textbf{Trust \& Quality Assurance}. Define: Hypothesis, Experiment, Success Metric, Timebox.
    \item Identify 2 major regulatory/legal risks in Cameroon for this model and mitigations.
\end{enumerate}
\tcblower
\textbf{Detailed Solution:}
\begin{enumerate}
    \item \textbf{Lean Canvas -- KmerJobs:}
    \begin{center}
    \begin{tabularx}{\linewidth}{|l|X|}\hline
    \rowcolor{popblueL} \textbf{Block} & \textbf{Content} \\ \hline
    \textbf{Problem} & 1. Clients: hard to find reliable artisans; fear of theft/overcharging; no reviews. 2. Artisans: no steady clients; underpaid; no digital presence; paid late. 3. Market: fragmented, cash-only, no dispute resolution. \\ \hline
    \textbf{Customer Segments} & Early adopters: (a) upper-middle households in Bonapriso/Akwa (high need, digital-savvy, pay premium); (b) SMEs (restaurants, hotels) needing maintenance; (c) artisans already using WhatsApp Business. \\ \hline
    \textbf{UVP} & ``Vetted artisans, fixed prices, and guaranteed work for Douala homes -- booked in 2 minutes via WhatsApp/App.'' \\ \hline
    \textbf{Solution (MVP)} & 1. Client: post job $\to$ get matched profiles $\to$ book. 2. Artisan app: accept job $\to$ navigate $\to$ complete $\to$ get paid (MoMo). 3. Admin dashboard: vetting, disputes, payouts. \\ \hline
    \textbf{Channels} & WhatsApp broadcast lists; artisan unions/associations; local radio spots; referral credit (``Refer a friend, get 2,000 FCFA''). \\ \hline
    \textbf{Revenue Streams} & 1. Transaction fee: 10\% per job. 2. Artisan Pro subscription: 5,000 FCFA/month for priority leads. 3. SME maintenance retainers (fixed monthly fee). \\ \hline
    \textbf{Cost Structure} & CAC (client) $\approx$ 3,000 FCFA; CAC (artisan) $\approx$ 1,500 FCFA via unions; COGS: MoMo fees, SMS, per-job insurance; Ops: small team + cloud hosting. \\ \hline
    \textbf{Key Metrics} & Supply: active artisans, fill rate $>80\%$. Demand: repeat rate $>40\%$ in 90 days, NPS $>50$. Unit economics: LTV/CAC $>3$, contribution margin $>30\%$. \\ \hline
    \textbf{Unfair Advantage} & 1. Partnership with artisan unions (access to vetted workers). 2. Proprietary Trust Score (ratings + union endorsement + ID verification + job history). 3. Per-job micro-insurance partnership. \\ \hline
    \end{tabularx}
    \end{center}
    \item \textbf{Market Sizing (Douala Year 1, illustrative assumptions):}
    \begin{align*}
    \text{TAM} &= 200,000 \text{ workers} \times 20 \text{ jobs/mo} \times 12 \times 15,000 = 720 \text{ Billion FCFA/yr (GMV)} \\
    \text{SAM} &\approx \text{digitally reachable share} \approx 500,000 \text{ jobs/yr} \times 15,000 = 7.5 \text{ Billion FCFA} \\
    \text{SOM (Y1: } \approx 0.5\% \text{ of SAM jobs)} &= 2,500 \text{ jobs} \times 15,000 = 37.5 \text{ Million FCFA GMV} \\
    \text{Platform Revenue (10\%)} &= 3.75 \text{ Million FCFA} \\
    \text{Subscriptions (100 artisans} \times 5,000 \times 12) &= 6 \text{ Million FCFA} \\
    \text{Total Year 1 Revenue} &\approx \mathbf{10 \text{ Million FCFA}}
    \end{align*}
    \emph{Note:} these figures are classroom assumptions for practising the method, not measured market data.
    \item \textbf{Concierge MVP Experiment:}
    \begin{itemize}
        \item \textbf{Hypothesis:} ``If we manually vet 20 artisans, provide fixed-price menus, and offer a `satisfaction or redo' guarantee, then 50 households will book at least 2 jobs/month via WhatsApp.''
        \item \textbf{Experiment:}
        \begin{enumerate}
            \item Recruit 20 artisans via the union; verify ID and skills; sign a service agreement.
            \item Create a fixed price menu (e.g., ``Unblock sink: 15,000 FCFA'').
            \item Launch a WhatsApp pilot group with 50 selected households.
            \item Client sends a job request; the founder manually matches the nearest available artisan, confirms the price, and sends the artisan's photo and trust score.
            \item Artisan executes; client confirms completion; admin triggers the MoMo payout and collects the 10\% fee.
            \item Post-job: client rates 1--5 stars with a comment.
        \end{enumerate}
        \item \textbf{Success Metrics (4 weeks):} $\ge 30$ unique clients book $\ge 1$ job; $\ge 60\%$ repeat booking; average rating $\ge 4.5/5$; zero safety incidents; positive unit economics (fee per job $>$ marginal cost of MoMo + SMS + insurance + admin time).
        \item \textbf{Timebox:} 4 weeks. \textbf{Decision:} metrics met $\to$ build the app; otherwise pivot (e.g., focus on SME maintenance contracts).
    \end{itemize}
    \item \textbf{Regulatory/Legal Risks (Cameroon):}
    \begin{enumerate}
        \item \textbf{Labour law / social security (CNPS):} If the platform controls prices, allocation, and discipline, artisans risk being \textbf{reclassified as employees} (CNPS contributions, Labour Code protections). \emph{Mitigation:} clear independent-contractor agreements; artisans set their own availability and may reject jobs; no exclusivity; obtain a legal opinion.
        \item \textbf{Consumer protection / liability:} As an intermediary, the platform may face claims for defective work, theft, or injury under general civil liability rules. \emph{Mitigation:} Terms of Service limiting liability to ``facilitation''; mandatory per-job micro-insurance; escrow payment released after client confirmation; KYC on both sides.
    \end{enumerate}
\end{enumerate}
\end{exoresolu}

\begin{methode}[Building a Minimum Viable Product (MVP) Roadmap]
\textbf{Objective:} Define the smallest feature set that delivers core value, enables learning, and de-risks the business.
\begin{enumerate}
    \item \textbf{Define the ``One Metric That Matters'' (OMTM) for the MVP:} e.g., ``50 artisans complete at least 3 jobs/week within 8 weeks of launch.''
    \item \textbf{User Story Mapping:}
    \begin{itemize}
        \item \textbf{Backbone:} User journey steps (Discover $\to$ Book $\to$ Execute $\to$ Pay $\to$ Review).
        \item \textbf{Slice 1 (Walking Skeleton):} End-to-end happy path only. No settings, no chat, no analytics, no admin panel (use SQL/Excel).
        \item \textbf{Slice 2 (MVP):} Add error handling, basic auth, notifications, artisan onboarding, payout automation.
        \item \textbf{Slice 3 (V1.1):} Ratings, disputes, search/filter, subscription billing.
    \end{itemize}
    \item \textbf{Tech Stack for Speed (Cameroon context):} a managed backend (e.g., Supabase: PostgreSQL + Auth + Realtime) to avoid DevOps; a single-codebase mobile framework (Flutter or React Native); MoMo/Orange Money APIs behind a unified payment wrapper; push + SMS notifications; simple static hosting for the web front-end.
    \item \textbf{Definition of Done (DoD):} Code $\to$ Unit test $\to$ Peer review $\to$ Staging deploy $\to$ Product-owner acceptance $\to$ Feature flag ON.
    \item \textbf{Launch Checklist:} Legal (Terms, Privacy), Security (OWASP Top 10 review), Monitoring (error tracking, logs), Support (WhatsApp Business), Analytics (privacy-friendly).
\end{enumerate}
\end{methode}

\begin{exoresolu}[MVP Scope \& Roadmap for ``KmerJobs'']
\textbf{Statement:} Based on the successful Concierge MVP (30 clients, 20 artisans, 60 jobs in 4 weeks, 4.6/5 rating), you are now funded (10M FCFA pre-seed) to build the KmerJobs App MVP. Team: 1 full-stack developer, 1 operations/community manager, 1 founder (product/sales). Timeline: 8 weeks.
\begin{enumerate}
    \item Create a User Story Map (Backbone + 3 Slices) for the \textbf{Client Side} only.
    \item Propose a database schema (tables and key columns) for: Users, Artisans, Jobs, Bookings, Reviews, Payouts.
    \item Design the ``Walking Skeleton'' Sprint Plan (2 weeks).
    \item Risk: ``Artisans don't adopt the app (they prefer WhatsApp calls).'' Propose 2 product features and 1 operational tactic to drive adoption.
\end{enumerate}
\tcblower
\textbf{Detailed Solution:}
\begin{enumerate}
    \item \textbf{User Story Map -- Client Side:}
    \begin{center}
    \begin{tabularx}{\linewidth}{|l|X|X|X|}\hline
    \rowcolor{popblueL} \textbf{Step} & \textbf{Walking Skeleton (W1--2)} & \textbf{MVP (W3--6)} & \textbf{V1.1 (W7--8+)} \\ \hline
    \textbf{1. Discover} & View artisan list (name, trade, rating, price); filter by trade only & Search by keyword; filter by distance, rating, availability; artisan profile page & Recommendations; favourites; categories \\ \hline
    \textbf{2. Book} & Select artisan $\to$ pick date/time $\to$ confirm $\to$ WhatsApp notify artisan & In-app booking with calendar; push/SMS notify; fixed price menu; add notes/photos & Recurring bookings; quote requests for custom jobs \\ \hline
    \textbf{3. Execute} & Artisan calls client; client confirms ``Done'' via WhatsApp to admin & In-app chat; artisan ``Start/Complete'' buttons; client confirms completion (OTP) & Live tracking; photo evidence \\ \hline
    \textbf{4. Pay} & Client pays cash; admin records manually in Excel & MoMo in-app payment (escrow); auto-split platform fee; instant payout to artisan; digital receipt & Wallet/top-up; corporate billing for SMEs \\ \hline
    \textbf{5. Review} & Admin collects rating via WhatsApp; manual entry & Post-job rating prompt (1--5 + tags + comment); public on profile & Dispute flow; verified badge; portfolio gallery \\ \hline
    \end{tabularx}
    \end{center}
    \item \textbf{Database Schema (tables and key columns):}
    \begin{center}
    \begin{tabularx}{\linewidth}{|l|X|}\hline
    \rowcolor{popblueL} \textbf{Table} & \textbf{Key Columns (and notes)} \\ \hline
    \texttt{profiles} & \texttt{id} (links to auth user), \texttt{role} (client/artisan/admin), \texttt{full\_name}, \texttt{phone} (unique), \texttt{created\_at}. Access rule: a user sees/edits only their own profile; admins see all. \\ \hline
    \texttt{artisans} & \texttt{id} (links to profiles), \texttt{trade}, \texttt{bio}, \texttt{skills}, \texttt{id\_verified}, \texttt{skill\_verified}, \texttt{trust\_score}, \texttt{is\_active}, \texttt{avg\_rating}, \texttt{jobs\_completed}, \texttt{momo\_number}, \texttt{location}. Public read of active artisans; artisan edits own record. \\ \hline
    \texttt{jobs} & \texttt{id}, \texttt{client\_id}, \texttt{trade}, \texttt{description}, \texttt{photos}, \texttt{status} (open/matched/in\_progress/completed/cancelled/disputed), \texttt{preferred\_date}, \texttt{location}, \texttt{address}, \texttt{fixed\_price}. Client manages own jobs. \\ \hline
    \texttt{bookings} & \texttt{id}, \texttt{job\_id} (unique), \texttt{artisan\_id}, \texttt{status} (pending/accepted/started/completed/confirmed/cancelled), \texttt{agreed\_price}, \texttt{platform\_fee}, \texttt{artisan\_net}, \texttt{client\_otp}, timestamps. Client and artisan see their own bookings. \\ \hline
    \texttt{reviews} & \texttt{id}, \texttt{booking\_id} (unique), \texttt{client\_id}, \texttt{artisan\_id}, \texttt{rating} (1--5), \texttt{tags}, \texttt{comment}, \texttt{artisan\_reply}. Public read; client writes own review. \\ \hline
    \texttt{payouts} & \texttt{id}, \texttt{booking\_id} (unique), \texttt{artisan\_id}, \texttt{amount}, \texttt{momo\_txn\_id}, \texttt{status} (pending/processing/paid/failed/reversed), \texttt{idempotency\_key} (unique), \texttt{paid\_at}. Artisan reads own payouts; only the payment service writes. \\ \hline
    \end{tabularx}
    \end{center}
    \item \textbf{Walking Skeleton Sprint Plan (2 weeks):}
    \begin{center}
    \begin{tabularx}{\linewidth}{|c|l|X|}\hline
    \rowcolor{popblueL} \textbf{Day} & \textbf{Goal} & \textbf{Definition of Done / Demo} \\ \hline
    Mon & Setup & Backend project, authentication (phone/email), app repository, CI to internal testing track; sign-up/login works. \\ \hline
    Tue & Profiles & Profiles table + access rules; artisan onboarding form (trade, bio, MoMo number, location); admin can verify in the database console. \\ \hline
    Wed & Jobs & Jobs table; client ``Post Job'' screen (trade, description, date, map location); job appears in the open list. \\ \hline
    Thu & Matching & Bookings table; admin ``Match'' action creates a booking; artisan receives a WhatsApp notification with job details and price. \\ \hline
    Fri & \textbf{Demo} & End-to-end live flow in 5 minutes: client posts $\to$ admin matches $\to$ artisan notified $\to$ job done $\to$ marked completed $\to$ manual payout recorded. \\ \hline
    Mon--Tue & Polish \& Deploy & Fix bugs, loading states, error messages; deploy to internal testers (5 pilot users run the flow unaided). \\ \hline
    Wed--Fri & Pilot Buffer & Ops onboards 10 real artisans (ID verified) and 10 real clients; collect feedback; target $\ge 80\%$ ``would use again''. \\ \hline
    \end{tabularx}
    \end{center}
    \item \textbf{Driving Artisan App Adoption (vs WhatsApp):}
    \begin{enumerate}
        \item \textbf{Feature: ``One-Tap Accept + Auto-Navigation''} -- Push notification: ``New Job: 15,000 FCFA, 2.3 km away''. One tap accepts and opens the navigation app with the client's address. \emph{Value:} saves time, reduces ``I'm lost'' calls.
        \item \textbf{Feature: ``Instant Payout on Confirmation''} -- When the client confirms completion (OTP), the MoMo payout is triggered instantly; the artisan receives an SMS confirmation within seconds. \emph{Value:} cash-flow certainty.
        \item \textbf{Operational Tactic: ``Launch Bonus + Weekly Leaderboard''} -- First 50 artisans: ``Complete 10 jobs in Month 1 $\to$ 10,000 FCFA bonus + `Founding Artisan' badge (priority leads)''. Weekly WhatsApp broadcast of the top 5 earners creates social proof and gamifies adoption.
    \end{enumerate}
\end{enumerate}
\end{exoresolu}

\begin{methode}[Pitching the IT Business Case (Executive Summary \& Financials)]
\textbf{Objective:} Communicate the opportunity, solution, traction, and ask clearly to investors/partners (e.g., angel networks, GICAM circles, diaspora investors, donor programmes).
\begin{enumerate}
    \item \textbf{Executive Summary (1 page / 30-second pitch):}
    \begin{itemize}
        \item \textbf{Hook:} a striking, honest statement of the pain (e.g., artisans lose income to idle time while households struggle to find a trusted plumber).
        \item \textbf{Problem/Solution:} one sentence each.
        \item \textbf{Traction:} concrete pilot results (jobs done, rating, revenue, incidents).
        \item \textbf{Market:} SAM/SOM figures with stated assumptions.
        \item \textbf{Business Model:} fees + subscriptions; unit economics (LTV, CAC, payback).
        \item \textbf{Team:} relevant experience of founders.
        \item \textbf{Ask:} amount, instrument, use of funds, 12-month milestone.
    \end{itemize}
    \item \textbf{Financial Model (3-Year Projection -- illustrative drivers):}
    \begin{center}
    \begin{tabularx}{\linewidth}{|X|c|c|c|}\hline
    \rowcolor{popblueL} \textbf{Metric} & \textbf{Year 1} & \textbf{Year 2} & \textbf{Year 3} \\ \hline
    Active Artisans & 500 & 2,500 & 10,000 \\ \hline
    Jobs/Month/Artisan & 4 & 6 & 8 \\ \hline
    Total Jobs/Year & 24,000 & 180,000 & 960,000 \\ \hline
    Avg Job Value (FCFA) & 15,000 & 18,000 & 20,000 \\ \hline
    \textbf{GMV (FCFA)} & \textbf{360M} & \textbf{3.24B} & \textbf{19.2B} \\ \hline
    Take Rate (blended) & 10\% & 10\% & 10\% \\ \hline
    Transaction Revenue & 36M & 324M & 1.92B \\ \hline
    Subscriptions & 3M & 30M & 120M \\ \hline
    \textbf{Total Revenue} & \textbf{39M} & \textbf{354M} & \textbf{2.04B} \\ \hline
    COGS (MoMo, SMS, insurance) & 12M & 97M & 550M \\ \hline
    \textbf{Gross Profit} & \textbf{27M} & \textbf{257M} & \textbf{1.49B} \\ \hline
    Opex (team, cloud, marketing) & 60M & 180M & 450M \\ \hline
    \textbf{EBITDA} & \textbf{-33M} & \textbf{+77M} & \textbf{+1.04B} \\ \hline
    \end{tabularx}
    \end{center}
    \emph{Reading:} Year 1 is investment phase (negative EBITDA); break-even is projected during Year 2; the model only works if artisan acquisition and retention targets are met.
    \item \textbf{Use of Funds (example 50M FCFA):} Product 40\%, Growth 30\%, Operations 20\%, Buffer 10\%.
    \item \textbf{Key Risks \& Mitigations Slide:} Regulatory (partner with licensed institutions), Trust (insurance + vetting), Supply (union partnership), Technology (managed, scalable backend), Competition (focus on the blue-collar niche).
\end{enumerate}
\end{methode}

\begin{exoresolu}[Investor Q\&A Simulation for KmerJobs]
\textbf{Statement:} You are pitching KmerJobs to a panel of Cameroonian angel investors. They ask the following tough questions. Provide concise, data-backed answers (max 3 sentences each).
\begin{enumerate}
    \item ``Why won't artisans just take the client's phone number and bypass your platform (disintermediation) after the first job?''
    \item ``Your Year 2 revenue assumes 2,500 active artisans. How do you recruit them at scale without burning cash? CAC of 1,500 FCFA seems low.''
    \item ``What happens if MTN changes their MoMo API terms or increases fees? You're 100\% dependent.''
    \item ``How do you handle a dispute where a client claims the artisan stole valuables? Your insurance only covers limited property damage.''
    \item ``You have no proprietary tech. A competitor could copy this in 3 months. What is your real moat?''
\end{enumerate}
\tcblower
\textbf{Detailed Solution:}
\begin{enumerate}
    \item \textbf{Disintermediation:} ``Our \textbf{Trust Score} and \textbf{Insurance} are platform-native: artisans lose their verified badge, priority leads, and insurance cover if they transact off-platform. Clients \emph{only} get the satisfaction guarantee and dispute resolution when booking via KmerJobs. The 10\% fee buys peace of mind -- cheaper than the risk of an unverified stranger in your home.''
    \item \textbf{Artisan Acquisition at Scale:} ``The 1,500 FCFA CAC is validated by our \textbf{union partnership}: we onboard cohorts of 50--100 artisans at union meetings at near-zero marginal cost, not one-by-one through ads. Year 2 adds a \textbf{referral programme} (an artisan earns a bonus for each referred artisan who completes 5 jobs) and WhatsApp-assisted onboarding. We target about 500 artisans per quarter through union partners in Douala, Yaoundé and Bafoussam.''
    \item \textbf{MoMo Dependency:} ``We abstract payments behind a \textbf{unified payment gateway} of our own, with \textbf{Orange Money} integrated in parallel as a fallback. For high-value SME contracts we are exploring direct bank channels. If MoMo fees rise materially, we transparently pass a small processing fee to the client and batch payouts to reduce per-transaction cost.''
    \item \textbf{Theft Dispute:} ``\textbf{KYC and background checks} on artisans are the first filter and deterrent. For disputes: \textbf{escrow} holds funds for 24 hours; the client files a claim; we freeze the payout and mediate using photos and chat logs; unresolved cases go to arbitration under our Terms of Service. Theft is a criminal matter: we require a police report and cooperate fully (artisan ID, job and location records). Platform liability is capped in the Terms of Service.''
    \item \textbf{Moat (Defensibility):} ``Technology is commoditisable; a \textbf{trust network} is not. Our moat is the \textbf{union partnership} (a vetted artisan base built through long relationship work), the \textbf{proprietary Trust Score} trained on completed-job data that a newcomer does not have, and \textbf{two-sided network effects}: more artisans $\to$ better fill rates $\to$ happier clients $\to$ more jobs $\to$ more artisans. Large e-commerce players focus on goods and white-collar services; we own the blue-collar niche.''
\end{enumerate}
\end{exoresolu}

\begin{aretenir}[Key Points of the Lesson]
\begin{itemize}
    \item A \cle{business need} is a gap between the current situation and the desired situation; always distinguish \cle{symptoms} from \cle{root causes}.
    \item Use \cle{PIECES} (Performance, Information, Economy, Control, Efficiency, Service) to classify needs, and document each need as ``The business needs to \emph{<action>} so that \emph{<benefit>}.''
    \item Elicit requirements by \cle{triangulation}: interviews + observation + document analysis; record them as \cle{user stories} with acceptance criteria; prioritise with \cle{MoSCoW}.
    \item Model processes \cle{AS-IS} and \cle{TO-BE} (BPMN); hunt waste with \cle{TIMWOODS}.
    \item Match needs to solution categories (Operational, Analytical, Collaborative, Customer-Facing, Infrastructure) and decide Buy vs Build with a \cle{weighted decision matrix}.
    \item Test viability with \cle{TELOS} feasibility (Technical, Economic, Legal, Operational, Schedule), respecting the Cameroonian context (ANTIC, BEAC/COBAC, OHADA, MINCOMMERCE, OAPI).
    \item Validate opportunities with the \cle{Lean Canvas}, \cle{TAM/SAM/SOM} sizing, and cheap experiments (smoke test, concierge MVP) before building the full product.
\end{itemize}
\end{aretenir}

\begin{attention}[Common Mistakes to Avoid]
\begin{itemize}
    \item Confusing a \emph{solution} with a \emph{need}: ``We need a website'' is a solution; the need is ``reach customers beyond our quarter''.
    \item Relying on a single interviewee: the loudest voice is not always representative.
    \item Quoting invented market statistics: in the exam, state your assumptions clearly and compute consistently from them.
    \item Forgetting non-functional requirements (offline use, USSD fallback, audit trail) which are decisive in the Cameroonian context.
    \item In weighted matrices, forgetting to multiply each score by its weight before summing.
\end{itemize}
\end{attention}

\begin{exercice}[Practice: From Need to Opportunity]
A cybercafé owner in Bamenda complains: ``Students queue for hours to print documents, my USB sticks spread viruses, and I lose money because customers claim they paid when my notebook says otherwise.''
\begin{enumerate}
    \item Identify three stakeholders and one pain point for each.
    \item Write two structured need statements (``The business needs to \ldots so that \ldots'').
    \item Classify each need with PIECES.
    \item Propose one IT solution per need and justify it briefly for the Cameroonian context (connectivity, payment).
    \item State one legal/regulatory point the owner should check before launching an online print-ordering service.
\end{enumerate}
\end{exercice}

\begin{corrige}[Exercise 1]
\begin{enumerate}
    \item \textbf{Stakeholders:} Owner (lost revenue, disputes over payments); Students/Customers (long queues, wasted time); Employees (stress, errors in the notebook); (bonus) Suppliers of paper/ink (unpredictable demand).
    \item \textbf{Need statements:} (a) The business needs to \emph{let customers submit print jobs online (WhatsApp/web form) and pay in advance by Mobile Money} so that \emph{queues disappear and every order is paid and traceable}. (b) The business needs to \emph{record every transaction automatically in a digital till/POS} so that \emph{payment disputes are eliminated and daily revenue is known exactly}.
    \item \textbf{PIECES:} (a) Service (customer experience) + Efficiency (workflow) + Economy (secured revenue). (b) Control (audit trail, fraud prevention) + Information (accurate records).
    \item \textbf{Solutions:} (a) A simple WhatsApp Business ordering flow or lightweight web form integrated with MTN MoMo/Orange Money payment -- works with low connectivity and the dominant local payment habits. (b) An affordable POS/till application (offline-capable, syncing when the network returns) producing daily reports.
    \item \textbf{Legal point:} Register/declare the e-commerce activity properly (trade register, MINCOMMERCE obligations) and respect personal-data protection duties (ANTIC) for customer files and phone numbers; keep OHADA-compliant sales records.
\end{enumerate}
\end{corrige}

\newpage

\coursec{Exercises}

\courssub{I check my knowledge}

\begin{exercice}[MCQ: Business needs and IT]
Which of the following is \textbf{not} a typical business need that can be solved with IT?
\begin{enumerate}
\item Automating payroll calculations for employees.
\item Replacing a broken office chair.
\item Tracking stock levels in real time.
\item Managing customer appointments digitally.
\end{enumerate}
\end{exercice}

\begin{exercice}[True/False with justification]
State whether each statement is True or False and justify your answer briefly.
\begin{enumerate}
\item A feasibility study only considers the financial cost of an IT project.
\item Mobile Money can be considered an IT solution for a business need related to secure payments.
\item The 5W method (Who, What, Where, When, Why) is used to write a clear problem statement.
\item A cybercafé owner who wants more customers has a business need that cannot be solved with IT.
\end{enumerate}
\end{exercice}

\begin{exercice}[Definitions]
Define the following terms in the context of identifying business needs for IT solutions:
\begin{enumerate}
\item Business need
\item Feasibility study
\item Problem statement (using the 5W method)
\item IT solution
\end{enumerate}
\end{exercice}

\begin{exercice}[Direct application: Matching needs to tools]
Match each business problem in Column A with the most appropriate IT tool from Column B. Write the pairs (e.g., 1--F).

\begin{center}
\begin{tabularx}{\linewidth}{|l|X|}\hline
\textbf{Column A: Business Problem} & \textbf{Column B: IT Tool} \\ \hline
1. Lost or disorganised paper records & A. Payroll software \\ \hline
2. Employees paid late due to manual calculations & B. Database management system \\ \hline
3. Unknown stock levels leading to over/under-stocking & C. Stock/inventory management software \\ \hline
4. Few customers and low visibility & D. Social media marketing / website \\ \hline
5. Cash handling risks and theft & E. Mobile Money / digital payment platform \\ \hline
6. Missed appointments and double bookings & F. Digital calendar / scheduling software \\ \hline
\end{tabularx}
\end{center}
\end{exercice}

\courssub{I practise}

\begin{exercice}[Scenario analysis: Madam Amina's restaurant]
Madam Amina runs a restaurant in Bamenda. She keeps orders in a paper notebook, often forgets orders, cannot know her daily profit, and customers complain about long waiting times.
\begin{enumerate}
\item Identify \textbf{three} distinct business needs from this scenario.
\item For each need, propose \textbf{one} suitable IT solution (name the type of software or tool).
\item Choose \textbf{one} of your proposed solutions and justify why it is appropriate, mentioning at least two benefits for Madam Amina.
\end{enumerate}
\end{exercice}

\begin{exercice}[Survey design: Cybercafé owner]
A student wants to identify IT-solvable business needs for a cybercafé owner in Yaoundé. Write \textbf{five} open-ended interview questions the student could ask. Each question must target a different potential area of need (e.g., sales, customer management, maintenance, marketing, finance).
\end{exercice}

\begin{exercice}[Problem statement: School canteen]
The school canteen at Government Bilingual High School X uses a paper ledger to record daily sales and expenses. The bursar spends hours every Friday reconciling accounts, errors are frequent, and the principal cannot get a quick overview of the canteen's financial health.
\begin{enumerate}
\item Write a complete \textbf{problem statement} using the \textbf{5W method} (Who, What, Where, When, Why).
\item Identify \textbf{two} IT solutions that could address this problem.
\end{enumerate}
\end{exercice}

\begin{exercice}[Feasibility dimensions]
List and explain the \textbf{four main dimensions} of a feasibility study for an IT project. For each dimension, give a concrete example relevant to a small business in Cameroon (e.g., a boutique, a farm cooperative, a transport agency).
\end{exercice}

\courssub{I prepare for the GCE}

\begin{exercice}[Structured essay: Feasibility study dimensions] (10 marks)
\begin{enumerate}
\item List the four dimensions of a feasibility study. (4 marks)
\item For each dimension, explain what it assesses and illustrate with a specific example from a Cameroonian context (e.g., a startup in Douala, a cooperative in the Northwest, an e-government service). (6 marks)
\end{enumerate}
\end{exercice}

\begin{exercice}[Case study: Nkambe Farmers' Cooperative] (15 marks)
The \emph{Nkambe Potato Farmers' Cooperative} (NPFC) in the Northwest Region produces high-quality potatoes. Currently, they sell to middlemen who come to the farm, often at low prices. The farmers want to sell directly to buyers in Douala (supermarkets, restaurants, households) to increase their income. They have smartphones but limited internet access in the village. The cooperative has a small office with one laptop and a printer in Nkambe town (30 minutes away).
\begin{enumerate}
\item Identify \textbf{three} business needs of the NPFC that can be solved with IT. (3 marks)
\item Propose \textbf{two} distinct IT solutions (e.g., a platform, an app, a system) that could help them reach buyers in Douala. For each solution, describe briefly how it would work. (6 marks)
\item Evaluate the \textbf{feasibility} of \textbf{one} of your proposed solutions using \textbf{three} dimensions of a feasibility study (Technical, Economic, Operational, or Legal/Social). (4 marks)
\item State \textbf{one} major limitation or risk for \textbf{each} of the two solutions you proposed. (2 marks)
\end{enumerate}
\end{exercice}

\begin{exercice}[Essay: From need to solution] (10 marks)
``Identifying a business need is only the first step; the real challenge is matching it to a feasible IT solution.''
\begin{enumerate}
\item Describe the \textbf{process} a student entrepreneur should follow to move from observing a business problem to proposing a justified IT solution. Structure your answer in clear steps. (6 marks)
\item Explain why it is important to involve the \textbf{end-users} (business owner, staff, customers) during this process. Give two reasons. (4 marks)
\end{enumerate}
\end{exercice}

\newpage

\coursec{Solutions to the exercises}

\begin{corrige}[Exercise 1 — MCQ: Business needs and IT]
\textbf{Correct answer: 2 — Replacing a broken office chair.}

\textbf{Justification:} A business need solvable with IT must involve the processing, storage, transmission or presentation of \cle{information}, or the automation of a task based on data.
\begin{itemize}
\item \textbf{Option 1 (payroll automation):} IT-solvable — payroll software computes salaries, taxes and payslips automatically from employee data.
\item \textbf{Option 3 (real-time stock tracking):} IT-solvable — inventory management software updates stock levels with each sale or delivery.
\item \textbf{Option 4 (digital appointment management):} IT-solvable — scheduling software manages bookings and sends reminders.
\item \textbf{Option 2 (broken office chair):} This is a \emph{physical/material} need (furniture repair or purchase). No information processing is involved, so IT cannot solve it.
\end{itemize}
\textbf{Examiner's tip:} In MCQs of this type, always ask: ``Does solving this problem require \emph{data} or \emph{software}?'' If the answer is no, it is not an IT-solvable need.
\end{corrige}

\begin{corrige}[Exercise 2 — True/False with justification]
\begin{enumerate}
\item \textbf{False.} A feasibility study considers \emph{several} dimensions, not only cost: \cle{technical} feasibility (do we have the technology and skills?), \cle{economic} feasibility (costs vs.\ benefits), \cle{operational} feasibility (will people actually use it?), and \cle{legal/social} feasibility (is it lawful and acceptable?). Finance is only one dimension.
\item \textbf{True.} Mobile Money (e.g., MTN MoMo, Orange Money) is an IT solution: it uses software platforms, databases and telecommunication networks to satisfy the business need of \emph{sending, receiving and securing payments} without handling physical cash.
\item \textbf{True.} The 5W method structures a problem statement by answering \emph{Who} is affected, \emph{What} the problem is, \emph{Where} it occurs, \emph{When} it occurs, and \emph{Why} it matters. This produces a clear, complete description of the need before any solution is proposed.
\item \textbf{False.} ``Wanting more customers'' is a \emph{marketing/visibility} need, which IT can address very well: a website, social media pages (Facebook, WhatsApp Business), online advertising or a customer loyalty application can all attract and retain customers.
\end{enumerate}
\end{corrige}

\begin{corrige}[Exercise 3 — Definitions]
\begin{enumerate}
\item \textbf{Business need:} A problem, difficulty or opportunity within a business activity (sales, stock, finance, communication, customer relations\ldots) whose resolution would improve the performance, profit or organisation of the business. In this chapter we focus on needs that can be satisfied through information technology.
\item \textbf{Feasibility study:} A systematic analysis carried out \emph{before} implementing an IT project, to determine whether the proposed solution is achievable and worthwhile. It examines four dimensions: technical, economic, operational and legal/social feasibility.
\item \textbf{Problem statement (5W method):} A short, precise written description of a business problem built by answering five questions: \emph{Who} is affected by the problem? \emph{What} exactly is the problem? \emph{Where} does it occur? \emph{When} does it occur (frequency, moments)? \emph{Why} is it a problem (consequences, losses)? It guides the search for an appropriate solution.
\item \textbf{IT solution:} A combination of hardware, software, data and procedures based on information technology, designed to satisfy an identified business need — for example a stock management application, a website, a Mobile Money payment service or a computerised payroll system.
\end{enumerate}
\end{corrige}

\begin{corrige}[Exercise 4 — Matching needs to tools]
\textbf{Correct pairs:} 1--B, \quad 2--A, \quad 3--C, \quad 4--D, \quad 5--E, \quad 6--F.

\textbf{Justifications:}
\begin{enumerate}
\item \textbf{1--B:} Lost or disorganised paper records are solved by a \emph{database management system} (DBMS), which stores records electronically, allows fast searching and prevents physical loss.
\item \textbf{2--A:} Late salaries caused by manual calculations are solved by \emph{payroll software}, which computes wages, deductions and payslips automatically and on time.
\item \textbf{3--C:} Unknown stock levels are solved by \emph{stock/inventory management software}, which updates quantities in real time and can trigger reorder alerts.
\item \textbf{4--D:} Low visibility is a marketing need, addressed by \emph{social media marketing or a website} to reach more customers.
\item \textbf{5--E:} Cash handling risks are reduced by a \emph{Mobile Money / digital payment platform}, which limits physical cash in the shop and keeps a transaction trail.
\item \textbf{6--F:} Missed appointments and double bookings are solved by a \emph{digital calendar / scheduling software}, which centralises bookings and sends reminders.
\end{enumerate}
\end{corrige}

\begin{corrige}[Exercise 5 — Scenario analysis: Madam Amina's restaurant]
\textbf{1. Three distinct business needs:}
\begin{itemize}
\item \textbf{Order management need:} orders are written in a paper notebook and often forgotten $\rightarrow$ need for reliable order recording and tracking.
\item \textbf{Financial management need:} she cannot know her daily profit $\rightarrow$ need for automatic recording of sales and expenses and profit calculation.
\item \textbf{Customer service need:} customers complain about long waiting times $\rightarrow$ need to speed up order taking, kitchen communication and billing.
\end{itemize}

\textbf{2. One suitable IT solution per need:}
\begin{itemize}
\item Order management $\rightarrow$ a \emph{Point of Sale (POS) / restaurant order management application} (on a tablet or smartphone) that records each order digitally and sends it to the kitchen.
\item Financial management $\rightarrow$ an \emph{accounting / sales tracking software} (even a well-designed spreadsheet) that totals daily sales, deducts expenses and computes profit automatically.
\item Customer service $\rightarrow$ a \emph{digital billing and order-display system} (POS with kitchen screen or printer) that reduces errors and waiting time; optionally a WhatsApp Business line for advance orders.
\end{itemize}

\textbf{3. Justification of one solution (example: the POS order management application):}

The POS application is appropriate because the root cause of both forgotten orders and long waiting times is the \emph{manual, paper-based} order process. Benefits include:
\begin{itemize}
\item \textbf{Accuracy:} every order is captured digitally and transmitted instantly to the kitchen, so no order is forgotten or misread.
\item \textbf{Speed:} waiters no longer walk to the kitchen with paper; bills are computed automatically, reducing customer waiting time and increasing table turnover (more customers served per day).
\item \textbf{Bonus:} most POS systems also produce daily sales reports, which partially solves the financial need as well.
\end{itemize}
\end{corrige}

\begin{corrige}[Exercise 6 — Survey design: Cybercafé owner]
Five open-ended questions, each targeting a different area of need:

\begin{enumerate}
\item \textbf{Sales/services:} ``Can you describe how you record the services each customer uses (printing, browsing time, scanning) and how you calculate what each customer must pay?''
\item \textbf{Customer management:} ``How do you keep in touch with your regular customers, and how do you know which services they ask for most often?''
\item \textbf{Maintenance:} ``What happens when one of your computers or your internet connection breaks down, and how do you keep track of repairs and their costs?''
\item \textbf{Marketing:} ``How do new customers in this neighbourhood find out about your cybercafé and the services you offer?''
\item \textbf{Finance:} ``How do you follow your daily income and expenses, and how do you know at the end of the month whether the business made a profit?''
\end{enumerate}

\textbf{Examiner's expectations:} questions must be \emph{open-ended} (they cannot be answered by yes/no), must concern the \emph{current situation and difficulties} (not directly propose a solution), and must each target a distinct functional area. Any equivalent formulation covering five different areas earns full credit.
\end{corrige}

\begin{corrige}[Exercise 7 — Problem statement: School canteen]
\textbf{1. Problem statement using the 5W method:}
\begin{itemize}
\item \textbf{Who:} The bursar of GBHS X (who reconciles the accounts), the canteen staff, and the principal (who needs financial reports).
\item \textbf{What:} Daily sales and expenses of the canteen are recorded by hand in a paper ledger; reconciling the accounts is slow and full of errors, and no quick financial overview is available.
\item \textbf{Where:} At the school canteen and the bursar's office of GBHS X.
\item \textbf{When:} Every day for data entry; every Friday, several hours are spent on weekly reconciliation; errors appear throughout the term.
\item \textbf{Why:} Because manual recording wastes the bursar's time, produces unreliable figures (risk of losses or undetected fraud), and prevents the principal from making timely decisions about the canteen.
\end{itemize}
\textbf{Condensed statement:} \emph{``At GBHS X, the bursar records canteen sales and expenses in a paper ledger; every Friday he spends hours reconciling error-prone accounts, and the principal cannot obtain a quick, reliable view of the canteen's financial health, which delays decisions and creates a risk of financial loss.''}

\textbf{2. Two IT solutions:}
\begin{itemize}
\item A \emph{computerised accounting / cash-book application} (or structured spreadsheet) that records each sale and expense, computes totals automatically and generates weekly reports.
\item A \emph{Point of Sale (POS) system with a small database} at the canteen counter, so that every sale is captured at the source and daily/weekly summaries and dashboards are produced automatically for the principal.
\end{itemize}
\end{corrige}

\begin{corrige}[Exercise 8 — Feasibility dimensions]
The four main dimensions of a feasibility study are:

\begin{enumerate}
\item \textbf{Technical feasibility:} assesses whether the required technology (hardware, software, network, electricity) and technical skills exist or can be obtained. \emph{Example:} a boutique in Bafoussam wants a stock management app — is there a reliable power supply, a smartphone or computer, and someone able to install and maintain the app?
\item \textbf{Economic (financial) feasibility:} compares the total cost of the project (equipment, licences, training, internet subscription, maintenance) with the expected benefits (time saved, losses avoided, extra revenue). \emph{Example:} a farm cooperative checks whether a $150\,000\ \mathrm{FCFA}$ investment in a laptop and accounting software will be recovered within a year through better price negotiation and fewer record-keeping errors.
\item \textbf{Operational feasibility:} assesses whether the solution fits the way the business actually works and whether the people concerned (owner, staff, customers) are willing and able to use it. \emph{Example:} a transport agency introduces digital ticket booking — will the ticket clerks, who are used to paper tickets, accept the change, and are they trained to use the system?
\item \textbf{Legal/social feasibility:} assesses whether the solution respects the law (data protection, taxation, licensing) and is socially acceptable to the community. \emph{Example:} a shop that collects customers' phone numbers for a loyalty programme must handle that personal data responsibly, in line with Cameroon's data protection framework (Law No. 2010/012 on cybersecurity and cybercriminality and ANTIC guidelines), and customers must trust the system.
\end{enumerate}
\end{corrige}

\begin{corrige}[Exercise 9 — Structured essay: Feasibility study dimensions (10 marks)]
\textbf{Indicative mark scheme and model answer.}

\textbf{(a) The four dimensions (4 marks — 1 mark each):}
\begin{enumerate}
\item Technical feasibility;
\item Economic (financial) feasibility;
\item Operational feasibility;
\item Legal/social feasibility.
\end{enumerate}

\textbf{(b) Explanation and Cameroonian illustration (6 marks — 1.5 mark per dimension: 1 for explanation, 0.5 for a valid example):}
\begin{itemize}
\item \textbf{Technical feasibility} assesses whether the technology and skills needed for the project exist or can be acquired: hardware, software, connectivity, electricity, maintenance capacity. \emph{Example:} a startup in Douala that wants to launch a delivery app must verify that its developers can build it and that customers have smartphones and internet access.
\item \textbf{Economic feasibility} assesses whether the project is financially worthwhile: it weighs development, equipment and running costs against expected savings and revenue. \emph{Example:} a cooperative in the Northwest compares the cost of a billing system with the money lost each season through unrecorded sales.
\item \textbf{Operational feasibility} assesses whether the solution will actually work in the real organisation: does it match existing procedures, and will staff and customers accept and use it? \emph{Example:} an e-government online tax declaration service will fail operationally if taxpayers are not trained or do not trust online declarations.
\item \textbf{Legal/social feasibility} assesses conformity with laws and regulations (data protection, cybercrime law, taxation) and social acceptability. \emph{Example:} any Cameroonian platform collecting citizens' personal data must comply with national data protection rules supervised by ANTIC.
\end{itemize}

\textbf{Examiner's expectations:} the four dimensions must be named precisely; definitions must show what is \emph{assessed}, not just restate the name; examples must be concrete and plausible. Vague answers such as ``technical = about technology'' without further explanation earn only partial credit.
\end{corrige}

\begin{corrige}[Exercise 10 — Case study: Nkambe Farmers' Cooperative (15 marks)]
\textbf{(a) Three IT-solvable business needs (3 marks — 1 mark each):}
\begin{enumerate}
\item \textbf{Market access / marketing need:} reaching buyers in Douala directly instead of depending on middlemen (visibility, advertising, customer contact).
\item \textbf{Communication need:} exchanging prices, quantities, photos and delivery arrangements with distant buyers quickly and cheaply.
\item \textbf{Sales and financial management need:} recording orders, payments and members' deliveries reliably (the office has a laptop that can host records); secure payment collection is also acceptable here.
\end{enumerate}

\textbf{(b) Two distinct IT solutions (6 marks — 3 marks each: 1 for the idea, 2 for a coherent description):}
\begin{itemize}
\item \textbf{Solution 1 — A WhatsApp Business / social media sales channel:} the cooperative creates a WhatsApp Business catalogue with photos, prices per bag and delivery conditions; a designated officer in Nkambe town (where connectivity is better) manages orders, confirms them by phone or message, and buyers pay via Mobile Money before or on delivery. Low cost, uses the smartphones the farmers already have.
\item \textbf{Solution 2 — A simple e-commerce / online marketplace presence:} the cooperative registers on an existing agricultural marketplace platform (or builds a simple website hosted from the office laptop) listing available quantities and prices; supermarkets and restaurants place orders online, and the cooperative organises grouped transport to Douala. Payments can be integrated through Mobile Money or bank transfer.
\end{itemize}

\textbf{(c) Feasibility evaluation of Solution 1 (WhatsApp Business channel) — 4 marks (any three dimensions, about 1.5 + 1.5 + 1):}
\begin{itemize}
\item \textbf{Technical:} feasible — farmers already own smartphones and WhatsApp is widely used; however, village connectivity is limited, so order management must be centralised in Nkambe town where the network is better.
\item \textbf{Economic:} highly feasible — the application is free; costs are limited to data bundles and the officer's time, while eliminating middlemen raises the price received per bag.
\item \textbf{Operational:} feasible with training — members must learn to send photos, update the catalogue and record orders; a clear procedure (one officer, fixed order days) makes adoption realistic.
\item \textbf{Legal/social:} Mobile Money payments leave a transaction trail, which is legally safer than cash; the cooperative should keep proper sales records for tax purposes.
\end{itemize}

\textbf{(d) One major limitation or risk per solution (2 marks — 1 mark each):}
\begin{itemize}
\item \textbf{Solution 1:} dependence on internet connectivity — poor network coverage in the village can delay order confirmation and make the channel unreliable during peak demand.
\item \textbf{Solution 2:} higher cost and complexity — building or maintaining a website/marketplace presence requires skills and money the cooperative may not have, and buyers may not trust an unknown platform (risk of fraud or low adoption).
\end{itemize}

\textbf{Examiner's expectations:} needs must come from the case text (middlemen, distance, smartphones, office equipment); solutions must be \emph{distinct} and described as working mechanisms, not just named; feasibility must explicitly use the dimension names; limitations must be realistic and specific to each solution.
\end{corrige}

\begin{corrige}[Exercise 11 — Essay: From need to solution (10 marks)]
\textbf{(a) The process from observed problem to justified IT solution (6 marks — 1 mark per well-explained step, any six coherent steps):}
\begin{enumerate}
\item \textbf{Observe and gather information:} visit the business, observe daily activities and interview the owner, staff and customers to detect difficulties (surveys, questionnaires, observation grids).
\item \textbf{Identify and list the business needs:} from the data collected, distinguish real needs from symptoms, and retain those that involve information processing (IT-solvable needs).
\item \textbf{Write a clear problem statement:} use the 5W method (Who, What, Where, When, Why) to describe the priority problem precisely.
\item \textbf{Analyse the causes and consequences:} understand why the problem exists (manual processes, lack of records, poor communication) and what it costs the business.
\item \textbf{Propose candidate IT solutions:} brainstorm possible tools (software, app, website, digital payment, database) that match the need, drawing on existing examples.
\item \textbf{Conduct a feasibility study:} evaluate each candidate solution on the technical, economic, operational and legal/social dimensions.
\item \textbf{Select and justify the best solution:} choose the most feasible option and defend the choice with arguments (benefits, cost, ease of adoption).
\item \textbf{Plan implementation and follow-up:} define equipment, training, budget and how success will be measured after deployment.
\end{enumerate}

\textbf{(b) Importance of involving end-users (4 marks — 2 marks per reason):}
\begin{itemize}
\item \textbf{Reason 1 — Accuracy of the need:} end-users (owner, staff, customers) live the problem daily; only they can describe what really happens, its frequency and its true causes. Without them, the entrepreneur risks solving the wrong problem or building a solution based on false assumptions.
\item \textbf{Reason 2 — Adoption and operational feasibility:} a solution imposed without consulting users is often rejected or misused. Involving users early ensures the tool matches their skills and habits, creates a sense of ownership, and guarantees that the solution will actually be used — which is exactly what operational feasibility measures.
\end{itemize}

\textbf{Examiner's expectations:} a structured, step-by-step answer (not a single paragraph); the 5W method and the four feasibility dimensions should appear explicitly; reasons in part (b) must be explained, not merely stated. A logical flow from observation to justification is the key marker of a top-band answer.
\end{corrige}

\newpage

\coursec{Summary sheet}

\begin{popfigure}
\begin{tikzpicture}[mindmap, grow cyclic, every node/.style={concept, text=white, font=\small\bfseries},
root concept/.append style={concept color=popdark, font=\bfseries},
level 1/.append style={level distance=3.6cm, sibling angle=60},
level 2/.append style={level distance=2.2cm, sibling angle=45, font=\scriptsize}]
\node[root concept]{Starting a Business: IT Needs}
  child[concept color=popblue]{ node{Businesses \& Needs}
    child{ node{Types of business} }
    child{ node{Needs: info, speed, reach} }
  }
  child[concept color=popgreen]{ node{Identifying Needs}
    child{ node{Observation} }
    child{ node{Interviews \& surveys} }
    child{ node{SWOT analysis} }
  }
  child[concept color=poporange]{ node{Matching to IT}
    child{ node{Software \& apps} }
    child{ node{Networks \& web} }
    child{ node{Mobile money} }
  }
  child[concept color=poppurple]{ node{Feasibility}
    child{ node{Technical} }
    child{ node{Economic} }
    child{ node{Legal (ANTIC)} }
  }
  child[concept color=poppink]{ node{Opportunities}
    child{ node{Gap in market} }
    child{ node{Business plan} }
    child{ node{Startup steps} }
  }
  child[concept color=popgold]{ node{GCE Skills}
    child{ node{Case studies} }
    child{ node{Justify choices} }
  };
\end{tikzpicture}
\legende{Mind map of the chapter ``Starting a business''.}
\end{popfigure}

\begin{bilan}
\textbf{Key definitions}
\begin{itemize}
\item \cle{Business}: an organisation that provides goods or services to customers in exchange for money, aiming at profit or a social goal.
\item \cle{Business need}: a problem to solve or an improvement to achieve (e.g.\ faster billing, better stock control, wider customer reach).
\item \cle{IT solution}: the use of hardware, software, networks or digital services to satisfy a business need.
\item \cle{Needs assessment}: the systematic process of collecting and analysing information to identify what a business truly requires.
\item \cle{SWOT analysis}: study of \textbf{S}trengths, \textbf{W}eaknesses (internal), \textbf{O}pportunities and \textbf{T}hreats (external).
\item \cle{Feasibility study}: evaluation of whether a proposed IT project is technically possible, economically profitable, legally acceptable and operationally realistic.
\item \cle{Business opportunity}: a gap in the market that an entrepreneur can fill with a product or service.
\end{itemize}

\textbf{Techniques for identifying business needs}
\begin{itemize}
\item \textbf{Observation}: watch daily activities (queues at a Douala shop, manual record books).
\item \textbf{Interviews and questionnaires}: ask owners, workers and customers about difficulties.
\item \textbf{Document analysis}: study existing records, invoices, registers.
\item \textbf{Brainstorming and SWOT}: group reflection to classify internal and external factors.
\end{itemize}

\textbf{Matching needs to IT solutions (typical pairs)}
\begin{itemize}
\item Slow manual billing $\rightarrow$ point-of-sale (POS) software.
\item Poor stock control $\rightarrow$ inventory / database management system.
\item Limited customer reach $\rightarrow$ website, social media, e-commerce.
\item Cash-handling risks $\rightarrow$ mobile money (MTN MoMo, Orange Money).
\item Communication costs $\rightarrow$ email, messaging, VoIP.
\item Data insecurity $\rightarrow$ backups, antivirus, access control.
\end{itemize}

\textbf{Feasibility criteria (TELOS)}
\begin{itemize}
\item \textbf{T}echnical: do the equipment, skills and infrastructure exist?
\item \textbf{E}conomic: do benefits exceed costs (hardware, licences, training)?
\item \textbf{L}egal: respect of laws (data protection, ANTIC guidelines, business registration).
\item \textbf{O}perational: will staff accept and use the system?
\item \textbf{S}chedule: can it be delivered in the required time?
\end{itemize}

\textbf{Method: from need to opportunity}
\begin{enumerate}
\item Observe a business and list its problems.
\item Confirm the real need through interviews or surveys.
\item Propose an IT solution matched to each need.
\item Check feasibility (TELOS) and estimate costs.
\item If viable, draft a simple business plan and launch.
\end{enumerate}

\textbf{Common mistakes}
\begin{itemize}
\item Confusing a \emph{need} (problem) with a \emph{solution} (tool): first state the problem.
\item Proposing technology without checking cost, electricity supply or staff skills.
\item Ignoring legal aspects (data protection, registration).
\item Writing SWOT factors in the wrong category (internal vs external).
\end{itemize}

\textbf{GCE tips}
\begin{itemize}
\item In case-study questions, always \textbf{justify} the IT solution by linking it explicitly to the stated need.
\item Use Cameroonian examples naturally: mobile money, cybercafés, e-government services.
\item Structure essays: definition $\rightarrow$ identification technique $\rightarrow$ matched solution $\rightarrow$ feasibility.
\item Define every technical term you use; examiners reward precise vocabulary.
\end{itemize}
\end{bilan}

\findecours

\end{document}
